\documentclass[11pt]{article}
\usepackage[margin=1.1in,headsep=18pt]{geometry}
\usepackage{amsmath, amsthm, amssymb}
\newcommand{\dgshorttitle}{By Enough, But Not Too Much}
\newcommand{\dgauthorshort}{B. R. Davis}
% Davis Geometric house style, shared by the papers in docs/papers/.
% Loaded with % Davis Geometric house style, shared by the papers in docs/papers/.
% Loaded with % Davis Geometric house style, shared by the papers in docs/papers/.
% Loaded with \input{../dgstyle} after amsmath and amsthm (do not load amssymb: newpxmath provides it).
% Needs (TeX Live): newpx, titlesec, caption, fancyhdr, microtype, xcolor, enumitem.
\usepackage[T1]{fontenc}
\usepackage{newpxtext}
\usepackage[slantedGreek]{newpxmath}
\renewcommand{\ttdefault}{lmtt}  % Latin Modern typewriter for code and paths
\linespread{1.06}
\usepackage[final]{microtype}
\usepackage{xcolor}
\definecolor{dgink}{HTML}{1F2328}
\definecolor{dgmuted}{HTML}{5A5F66}
\definecolor{dgrule}{HTML}{C9CDD2}
\definecolor{dgchord}{HTML}{2A78D6}
\definecolor{dgarc}{HTML}{1BAF7A}
\definecolor{dgtan}{HTML}{EB6834}
\definecolor{dgblend}{HTML}{6250D6}
% headings: numbered small caps, quiet rules of space rather than ink
\usepackage{titlesec}
\titleformat{\section}{\normalfont\large\scshape\color{dgink}}{\thesection}{1em}{}
\titleformat{\subsection}{\normalfont\normalsize\itshape\color{dgink}}{\thesubsection}{0.8em}{}
\titleformat{\subsubsection}[runin]{\normalfont\normalsize\itshape}{\thesubsubsection}{0.6em}{}[.]
\titleformat{\paragraph}[runin]{\normalfont\normalsize\scshape}{}{0pt}{}  % sources end paragraph titles with their own stop
\titlespacing*{\section}{0pt}{2.4ex plus 0.6ex minus 0.3ex}{1.1ex plus 0.2ex}
\titlespacing*{\subsection}{0pt}{1.8ex plus 0.5ex minus 0.2ex}{0.7ex plus 0.2ex}
\titlespacing*{\paragraph}{0pt}{1.2ex plus 0.4ex minus 0.2ex}{0.6em}
% captions: small, label in small caps, ragged
\usepackage[font=small,labelfont=sc,labelsep=period,justification=justified,margin=4pt]{caption}
% running heads
\usepackage{fancyhdr}
\pagestyle{fancy}
\fancyhf{}
\renewcommand{\headrulewidth}{0pt}
\fancyhead[L]{\small\scshape\color{dgmuted}\dgauthorshort}
\fancyhead[R]{\small\itshape\color{dgmuted}\dgshorttitle}
\fancyfoot[C]{\small\color{dgmuted}\thepage}
\fancypagestyle{plain}{\fancyhf{}\fancyfoot[C]{\small\color{dgmuted}\thepage}}
\setlength{\headheight}{14pt}
% title block
\makeatletter
\renewcommand{\maketitle}{%
  \thispagestyle{plain}%
  \begin{center}
    {\LARGE\color{dgink}\@title\par}\vspace{1.1em}
    {\large\scshape\@author\par}\vspace{0.5em}
    {\small\color{dgmuted}\@date\par}
  \end{center}\vspace{0.6em}}
\renewenvironment{abstract}{%
  \begin{list}{}{\setlength{\leftmargin}{2.2em}\setlength{\rightmargin}{2.2em}}\item\relax
  \small\noindent{\scshape Abstract.}\enspace}{\end{list}\vspace{0.4em}}
\makeatother
\usepackage{enumitem}
\setlist{itemsep=0.25ex, topsep=0.6ex}
% "In plain words": a quiet ruled aside
\newenvironment{plain}{%
  \par\smallskip\begin{list}{}{\setlength{\leftmargin}{1.4em}\setlength{\rightmargin}{1.4em}}\item\relax
  \small\color{dgink}\noindent{\scshape In plain words.}\enspace}{\end{list}\smallskip}
 after amsmath and amsthm (do not load amssymb: newpxmath provides it).
% Needs (TeX Live): newpx, titlesec, caption, fancyhdr, microtype, xcolor, enumitem.
\usepackage[T1]{fontenc}
\usepackage{newpxtext}
\usepackage[slantedGreek]{newpxmath}
\renewcommand{\ttdefault}{lmtt}  % Latin Modern typewriter for code and paths
\linespread{1.06}
\usepackage[final]{microtype}
\usepackage{xcolor}
\definecolor{dgink}{HTML}{1F2328}
\definecolor{dgmuted}{HTML}{5A5F66}
\definecolor{dgrule}{HTML}{C9CDD2}
\definecolor{dgchord}{HTML}{2A78D6}
\definecolor{dgarc}{HTML}{1BAF7A}
\definecolor{dgtan}{HTML}{EB6834}
\definecolor{dgblend}{HTML}{6250D6}
% headings: numbered small caps, quiet rules of space rather than ink
\usepackage{titlesec}
\titleformat{\section}{\normalfont\large\scshape\color{dgink}}{\thesection}{1em}{}
\titleformat{\subsection}{\normalfont\normalsize\itshape\color{dgink}}{\thesubsection}{0.8em}{}
\titleformat{\subsubsection}[runin]{\normalfont\normalsize\itshape}{\thesubsubsection}{0.6em}{}[.]
\titleformat{\paragraph}[runin]{\normalfont\normalsize\scshape}{}{0pt}{}  % sources end paragraph titles with their own stop
\titlespacing*{\section}{0pt}{2.4ex plus 0.6ex minus 0.3ex}{1.1ex plus 0.2ex}
\titlespacing*{\subsection}{0pt}{1.8ex plus 0.5ex minus 0.2ex}{0.7ex plus 0.2ex}
\titlespacing*{\paragraph}{0pt}{1.2ex plus 0.4ex minus 0.2ex}{0.6em}
% captions: small, label in small caps, ragged
\usepackage[font=small,labelfont=sc,labelsep=period,justification=justified,margin=4pt]{caption}
% running heads
\usepackage{fancyhdr}
\pagestyle{fancy}
\fancyhf{}
\renewcommand{\headrulewidth}{0pt}
\fancyhead[L]{\small\scshape\color{dgmuted}\dgauthorshort}
\fancyhead[R]{\small\itshape\color{dgmuted}\dgshorttitle}
\fancyfoot[C]{\small\color{dgmuted}\thepage}
\fancypagestyle{plain}{\fancyhf{}\fancyfoot[C]{\small\color{dgmuted}\thepage}}
\setlength{\headheight}{14pt}
% title block
\makeatletter
\renewcommand{\maketitle}{%
  \thispagestyle{plain}%
  \begin{center}
    {\LARGE\color{dgink}\@title\par}\vspace{1.1em}
    {\large\scshape\@author\par}\vspace{0.5em}
    {\small\color{dgmuted}\@date\par}
  \end{center}\vspace{0.6em}}
\renewenvironment{abstract}{%
  \begin{list}{}{\setlength{\leftmargin}{2.2em}\setlength{\rightmargin}{2.2em}}\item\relax
  \small\noindent{\scshape Abstract.}\enspace}{\end{list}\vspace{0.4em}}
\makeatother
\usepackage{enumitem}
\setlist{itemsep=0.25ex, topsep=0.6ex}
% "In plain words": a quiet ruled aside
\newenvironment{plain}{%
  \par\smallskip\begin{list}{}{\setlength{\leftmargin}{1.4em}\setlength{\rightmargin}{1.4em}}\item\relax
  \small\color{dgink}\noindent{\scshape In plain words.}\enspace}{\end{list}\smallskip}
 after amsmath and amsthm (do not load amssymb: newpxmath provides it).
% Needs (TeX Live): newpx, titlesec, caption, fancyhdr, microtype, xcolor, enumitem.
\usepackage[T1]{fontenc}
\usepackage{newpxtext}
\usepackage[slantedGreek]{newpxmath}
\renewcommand{\ttdefault}{lmtt}  % Latin Modern typewriter for code and paths
\linespread{1.06}
\usepackage[final]{microtype}
\usepackage{xcolor}
\definecolor{dgink}{HTML}{1F2328}
\definecolor{dgmuted}{HTML}{5A5F66}
\definecolor{dgrule}{HTML}{C9CDD2}
\definecolor{dgchord}{HTML}{2A78D6}
\definecolor{dgarc}{HTML}{1BAF7A}
\definecolor{dgtan}{HTML}{EB6834}
\definecolor{dgblend}{HTML}{6250D6}
% headings: numbered small caps, quiet rules of space rather than ink
\usepackage{titlesec}
\titleformat{\section}{\normalfont\large\scshape\color{dgink}}{\thesection}{1em}{}
\titleformat{\subsection}{\normalfont\normalsize\itshape\color{dgink}}{\thesubsection}{0.8em}{}
\titleformat{\subsubsection}[runin]{\normalfont\normalsize\itshape}{\thesubsubsection}{0.6em}{}[.]
\titleformat{\paragraph}[runin]{\normalfont\normalsize\scshape}{}{0pt}{}  % sources end paragraph titles with their own stop
\titlespacing*{\section}{0pt}{2.4ex plus 0.6ex minus 0.3ex}{1.1ex plus 0.2ex}
\titlespacing*{\subsection}{0pt}{1.8ex plus 0.5ex minus 0.2ex}{0.7ex plus 0.2ex}
\titlespacing*{\paragraph}{0pt}{1.2ex plus 0.4ex minus 0.2ex}{0.6em}
% captions: small, label in small caps, ragged
\usepackage[font=small,labelfont=sc,labelsep=period,justification=justified,margin=4pt]{caption}
% running heads
\usepackage{fancyhdr}
\pagestyle{fancy}
\fancyhf{}
\renewcommand{\headrulewidth}{0pt}
\fancyhead[L]{\small\scshape\color{dgmuted}\dgauthorshort}
\fancyhead[R]{\small\itshape\color{dgmuted}\dgshorttitle}
\fancyfoot[C]{\small\color{dgmuted}\thepage}
\fancypagestyle{plain}{\fancyhf{}\fancyfoot[C]{\small\color{dgmuted}\thepage}}
\setlength{\headheight}{14pt}
% title block
\makeatletter
\renewcommand{\maketitle}{%
  \thispagestyle{plain}%
  \begin{center}
    {\LARGE\color{dgink}\@title\par}\vspace{1.1em}
    {\large\scshape\@author\par}\vspace{0.5em}
    {\small\color{dgmuted}\@date\par}
  \end{center}\vspace{0.6em}}
\renewenvironment{abstract}{%
  \begin{list}{}{\setlength{\leftmargin}{2.2em}\setlength{\rightmargin}{2.2em}}\item\relax
  \small\noindent{\scshape Abstract.}\enspace}{\end{list}\vspace{0.4em}}
\makeatother
\usepackage{enumitem}
\setlist{itemsep=0.25ex, topsep=0.6ex}
% "In plain words": a quiet ruled aside
\newenvironment{plain}{%
  \par\smallskip\begin{list}{}{\setlength{\leftmargin}{1.4em}\setlength{\rightmargin}{1.4em}}\item\relax
  \small\color{dgink}\noindent{\scshape In plain words.}\enspace}{\end{list}\smallskip}

\usepackage{graphicx}
\usepackage{booktabs}
\usepackage{url}
\usepackage[hidelinks]{hyperref}

\newtheorem{theorem}{Theorem}[section]
\newtheorem{proposition}[theorem]{Proposition}
\newtheorem{lemma}[theorem]{Lemma}
\newtheorem{corollary}[theorem]{Corollary}
\theoremstyle{definition}
\newtheorem{definition}[theorem]{Definition}
\theoremstyle{remark}
\newtheorem{remark}[theorem]{Remark}
\newtheorem{conjecture}[theorem]{Conjecture}

\newcommand{\su}{\mathfrak{su}(2)}
\newcommand{\norm}[1]{\lVert #1\rVert}
\newcommand{\RR}{\mathbb{R}}

\title{By Enough, But Not Too Much\\[4pt] \large Where the Twist Is the Floor of the Yang--Mills Hessian on Berger Spheres}
\author{Bee Rosa Davis\\[1pt] {\normalsize\upshape\itshape Davis Geometric}\\[1pt]
{\footnotesize\upshape ORCID \href{https://orcid.org/0009-0009-8034-4308}{0009-0009-8034-4308}}}
\date{November 2026 (draft)}

\begin{document}
\maketitle

\begin{abstract}
On the Berger spheres $(t,1,1)$, every left-invariant $SU(2)$ field that is symmetric under the rotation about the
long circle is, up to a constant gauge transformation, $M_0=x\,\tau_0$, $M_1=y\,\tau_1$, $M_2=y\,\tau_2$: a strength
$x$ along the long circle and a strength $y$ across it. At every such field the twist
$f\,(\hat\theta^1+i\hat\theta^2)\otimes(\tau_1+i\tau_2)$ is an exact gauge-orthogonal eigenvector of the
Yang--Mills Hessian in the sector of spin $j/2$, with eigenvalue $(x+(j+2)/\sqrt t)^2+2\sqrt t\,x-y^2$. The set of $x$
where it is the floor of every sector is bounded on both sides, an interval $[x_{\mathrm{lo}},x^*]$: above by the
edge, below by a state of the sector $j=0$. For $x>0$ it is never the floor of every sector, and the region is bounded
below; both are proved. Where the charge-$0$ block of the sector $j=4$ competes, the edge computed in the sectors
$j\le12$ is a root of an explicit factor $P_2$ of degrees $7$, $8$ and $16$ in $x$, $y$ and $\sqrt t$, whose
restriction to the Berger path contains the polynomial of the threshold $t_A$ of the Floor paper. These polynomial
identities are also checked in Lean~4. On the domain $x\in[-3,-1.5]$, $y\in[0.5,3]$, $t\in[3,10]$, interval
certificates prove that the twist is the floor of every sector, with multiplicity two. The all-sector edge is at most
$\min(x^*_{12},0)$, where $x^*_{12}$ is the edge computed in the sectors $j\le12$; equality is conjectured for
$t\ge3$ and is refuted near $t=2.01$. Three simpler readings, each written down before its data, failed: that the
rotation symmetry makes the twist the floor, that a field strength aligned with its spin does, and that alignment
together with the curvature term singling it out does. Three further tests were pre-registered. At $48$ points never computed before, the twist's
region is one interval and no sector $13\le j\le79$ undercuts it, while the edge is the $P_2$ root at $34$ of $38$ points
and set by other blocks at $4$. In an $SU(3)$ family whose fibre field is free, it is the floor at none of $20$ shapes.
The third test is running.
\end{abstract}

\noindent\textbf{Keywords.} Yang--Mills Hessian; Berger sphere; Peter--Weyl decomposition; Weitzenb\"ock identity;
spectral gap; level crossing; computer-assisted proof; formal verification.

\noindent\textbf{MSC 2020.} 58J50, 81T13, 58E15, 53C30, 65G30, 03B35.

% ---------------------------------------------------------------------------------------------------
\section{Introduction}\label{sec:intro}

The second variation of the Yang--Mills energy says which directions at a gauge field lead downhill and which lead
up. In~\cite{davisfloor2026} I computed it at the \emph{Berger field}, the Levi-Civita connection of the Berger
metric $(t,1,1)$ on $S^3=SU(2)$ read as an $SU(2)$ connection. In every Peter--Weyl sector $j$ (spin $j/2$) the
field $b=f\,(\hat\theta^1+i\hat\theta^2)\otimes(\tau_1+i\tau_2)$, with $f$ the top weight along the long axis, is an
exact eigenvector of the Hessian, orthogonal to every gauge direction, with eigenvalue
\[
  \lambda_t(j)=\frac{j(j+2)-2(t-3)j-2(t-2)(t+4)}{t}.
\]
I call this field Helike (``she who turns''); below it is \emph{the twist}. The paper~\cite{davisfloor2026} proved
that $\lambda_t(j)$ is the lowest eigenvalue of every sector for every $t\in[2.75,10]$. Below that range it found, in
double precision, that the floor is lost below $t^*\approx2.6980$ for the gauge-orthogonal spectrum and below the
algebraic threshold $t_A=2.7311222782\ldots$ for $H+GG^*$.

The Berger field is one point of a larger family. This paper asks where in that family the twist is the floor, and
what sets the edge of that region. Stability of Yang--Mills fields on spheres and homogeneous spaces is
classical~\cite{bl1981,laquer1984,ohnitapan1990,kot1986}, results for minimal-energy Yang--Mills connections on
three-manifolds of nonnegative Ricci curvature are in~\cite{stern2010}, and a recent quantitative upper bound for the
lowest Hessian eigenvalue of abelian Yang--Mills--Higgs critical points on the round $S^3$ is~\cite{jin2026}. Those
results concern critical points. The fields here are in general not critical points. On Berger spheres, spectra are
known for the scalar Laplacian~\cite{lauret2019}, the Dirac operator~\cite{bar1992,klingschueth2022} and the Hodge
Laplacian on $1$-forms~\cite{henkellauret2026}; B\"ar's computation~\cite{bar1992} uses Frobenius reciprocity, the
same sector decomposition used here.

\paragraph{Main results.} Section~\ref{sec:setting} fixes the operator and the family $M_0=x\tau_0$, $M_1=y\tau_1$,
$M_2=y\tau_2$.
\begin{itemize}
  \item \emph{The family is general} up to constant gauge, within the class of left-invariant fields symmetric under
    the rotation about the long circle (Proposition~\ref{prop:general}).
  \item \emph{Charge blocks} (Lemma~\ref{lem:blocks}): a conserved charge splits every sector into blocks of at most
    nine states, with rows in closed form in $(x,y,t)$.
  \item \emph{The twist across the family} (Section~\ref{sec:twist}): its eigenvalue is
    $\lambda=(x+(j+2)/\sqrt t)^2+2\sqrt t\,x-y^2$ (Proposition~\ref{prop:closed}). The cost of lifting it from sector
    $0$ to sector $j\ge1$ carries a discount $2xj/\sqrt t$, negative if and only if $x<0$ (Corollary~\ref{cor:lift}).
    Its anatomy is exact: fibre momentum squared, a lowest-Landau-level transverse term, and a curvature term
    (Proposition~\ref{prop:anatomy}).
  \item \emph{The edge polynomial} (Section~\ref{sec:edge}): the determinant numerator for the charge-$0$ block of
    sector $4$ factors as $-16384\,P_1P_2$ (Proposition~\ref{prop:edgepoly}), and on the Berger path
    $P_2=-32(s-1)(s+1)Q_{16}(s)$ with $Q_{16}$ the polynomial of $t_A$ (Proposition~\ref{prop:bergerrestriction}).
    Both identities, and a bracket for the crossing, are checked in Lean~4: sixteen theorems in all
    (Section~\ref{subsec:lean}).
  \item \emph{For $x>0$ the twist is never the floor of every sector}; high sectors undercut it
    (Section~\ref{sec:interval}).
  \item \emph{The region is bounded below} by the sector-$0$ state $\hat\theta^0\otimes\tau_0$, so in $x$ it lies in an
    interval $[x_{\mathrm{lo}},x^*]$ (Section~\ref{sec:interval}).
  \item \emph{An exact ``if and only if''} for the sectors $j\le12$ on two slices ($t=3$ with $y=\sqrt3$ and with
    $y=3/2$), conditional on the closed-form blocks being the blocks of $H+GG^*$ (Section~\ref{sec:interval}).
  \item \emph{Large-sector slab proofs} for every sector $j\ge40$, on boxes reaching to within about $0.1$ of the edge
    ($0.02$ in one slab), and \emph{the floor on a domain}
    $D=\{x\in[-3,-1.5],\,y\in[0.5,3],\,t\in[3,10]\}$, every sector, by interval certificates
    (Section~\ref{sec:domain}).
\end{itemize}
The theorem on $D$ also appears in the forthcoming version~3 of~\cite{davisfloor2026}.
% Bee: confirm this sentence about version 3.

\paragraph{What is computed.} In double precision: tables of the edge in the sectors $j\le12$, and a map of the
eight charge blocks that set it (Section~\ref{subsec:competitors}). Along the Berger path the $H+GG^*$ edge is crossed
at $t_A$ to $9\times10^{-16}$ and the gauge-orthogonal edge at $t^*$ to $1.1\times10^{-12}$
(Section~\ref{subsec:berger}). Every positive edge found within $j\le12$ is an artefact of the cut-off: at all $58$
such grid points a higher sector undercuts the twist ($54$ in the map's tail run, the remaining four in an
independent check; Section~\ref{sec:interval}).

\paragraph{What failed.} Three readings were tested, each written down before its data, and each failed: that the
rotation symmetry makes the twist the floor, that a field strength aligned with the twist's spin does, and that
alignment together with the curvature term singling the twist out does (Section~\ref{sec:decide}). A model of the
edge as $c_1/t+c_2/t^2$ failed both of its pre-registered predictions. The conjecture that the all-sector edge equals
$\min(x^*_{12},0)$ is refuted near $t=2.01$ and is kept only for $t\ge3$. The statement that the edge is the largest
real root of $P_2$ is false: $P_2$ has further real roots where a higher eigenvalue of the competitor meets the twist.
Three further tests were pre-registered (Section~\ref{sec:prereg}). At held-out points the edge clause failed at $4$ of
$38$ points, where other blocks reach the twist first. In an $SU(3)$ family with a free fibre field the twist is the
floor nowhere, so the reading of Section~\ref{sec:scope} does not carry over to that family. The test on general
circle-symmetric fields is running.

\paragraph{By enough, but not too much.} The field along the long circle must turn against the twist by enough that
the discount $2xj/\sqrt t$ carries her below her competitor, but not by so much that the sector-$0$ state
$\hat\theta^0\otimes\tau_0$ undercuts her: $x_{\mathrm{lo}}\le x\le x^*$.

The thresholds $t^*$ and $t_A$ of~\cite{davisfloor2026} are therefore not special to the Berger field. They are where
the Berger path meets the edge. The literature I checked (the chromomagnetic
instability~\cite{nielsenolesen1978,savvidy1977}, homogeneous backgrounds on the round
$S^3$~\cite{avramidicollopy2012,klc2021}) did not contain such an edge; the nearest is the negative modes of covariantly
constant backgrounds on the round sphere~\cite{avramidicollopy2012}.

Section~\ref{sec:setting} sets up the operator, the family and the charge blocks. Section~\ref{sec:twist} gives the
twist's eigenvalue and its anatomy. Section~\ref{sec:edge} derives the edge, its polynomial, the Berger crossings and the
Lean check. Section~\ref{sec:interval} proves that the region is bounded on both sides and that $x>0$ loses.
Section~\ref{sec:domain} proves the floor on a domain. Section~\ref{sec:decide} reports what decides the edge and what
does not. Section~\ref{sec:prereg} describes the pre-registered tests and Section~\ref{sec:scope} the scope and open
problems; the closing statement on availability lists the code and data.

% ---------------------------------------------------------------------------------------------------
\section{Setting: the operator, the circle-symmetric family, charge blocks}\label{sec:setting}

\subsection{Conventions}\label{subsec:conventions}

The conventions are those of~\cite{davisfloor2026}, stated compactly. $S^3=SU(2)$ carries the left-invariant frame
$e_0,e_1,e_2$ with $[e_a,e_b]=2\epsilon_{abc}e_c$. The Berger metric $(t,1,1)$ is diagonal in it, with $e_0$ the long
axis; the parameter $t>0$ is \emph{the stretch}, and the orbits of $e_0$ are \emph{the long circle} (the fibre).
$\hat\theta^i$ is the orthonormal coframe. A field is $A=\sum_iM_i\hat\theta^i$ with $M_i$ constant in $\su$, which
has the basis $\tau_0,\tau_1,\tau_2$ with the same brackets, $\tau_0$ along the long axis. Its curvature is
$F_{ij}=[M_i,M_j]-\sum_kc_{ij}^{\,k}M_k$.

The second variation is $Q(b)=\sum_{i<j}\norm{D_{ij}b}^2+2\langle F_{ij},[b_i,b_j]\rangle$, and $H$ is its operator
against the $L^2$ mass $\sum_i\norm{b_i}^2$. The gauge map is $G\xi=(\hat e_i\xi+[M_i,\xi])_i$. The fields studied are
in general not Yang--Mills critical points, so gauge directions are not zero modes; the \emph{gauge-orthogonal
spectrum} is the spectrum of $PHP$ on $\ker G^*$. In the Peter--Weyl sector $j\ge0$ (spin $j/2$), $b_i\in V_j\otimes\su$
and $H$ is a $9(j+1)\times9(j+1)$ Hermitian matrix. The Ricci tensor of $(t,1,1)$ is $\mathrm{diag}(2t,4-2t,4-2t)$ in
the orthonormal frame.

The operator $H+GG^*$ has the Weitzenb\"ock form
\begin{equation}\label{eq:weitz}
  H+GG^*=\sum_iX_i^*X_i+R_0,\qquad X_i=\rho_i/\sqrt{g_i}+\Omega_i,\qquad R_0=\mathrm{Ric}\otimes I-2[F,\cdot\,],
\end{equation}
with $([F,\cdot\,]b)_i=\sum_j[F_{ij},b_j]$: the Ricci curvature plus the field strength with the factor $2$ of the
Hessian, as in the Bochner--Weitzenb\"ock formula of~\cite{bl1981} (an accessible statement is~\cite{hanwen2026}). The
split of $R_0$ is checked on the operator in Remark~\ref{rem:split}. Here $\rho_0,\rho_1,\rho_2=-2iJ_z,-2iJ_x,-2iJ_y$
act on $V_j$, and $\Omega_i$ is the connection part on the frame and Lie indices.

The \emph{floor} of a sector is its lowest gauge-orthogonal eigenvalue, or, where stated, its lowest eigenvalue of
$H+GG^*$. The \emph{edge} is the level-crossing locus where the twist's eigenvalue meets the lowest eigenvalue of
another charge block, \emph{the competitor}.

\subsection{The circle-symmetric family}\label{subsec:family}

The family is
\[
  M_0=x\,\tau_0,\qquad M_1=y\,\tau_1,\qquad M_2=y\,\tau_2 ,
\]
a strength $x$ along the long circle and a strength $y$ across it. Three members recur. The Berger field is
$x=(2-t)/\sqrt t$, $y=\sqrt t$; I call the curve it traces \emph{the Berger path}. The multiples $A=s\theta$ of the
Maurer--Cartan form are $x=2s/\sqrt t$, $y=2s$. The squashed meron lies in the family too, at $x>0$ (its coordinates
are computed numerically). The field strength of a member is
\begin{equation}\label{eq:B}
  F_{12}=B\,\tau_0,\quad B=y^2-2\sqrt t\,x,\qquad F_{01}=(xy-2y/\sqrt t)\,\tau_2,\quad F_{02}=-(xy-2y/\sqrt t)\,\tau_1 ,
\end{equation}
so $B$ is the field along the fibre and $F_{01},F_{02}$, of size $|y(x-2/\sqrt t)|$, lie across it. On the Berger path
$B=3t-4$; on $A=s\theta$, $B=4s(s-1)$ (\path{scripts/ricci_su2_twist_anatomy.py}, receipt \path{twist_anatomy.json}).

\begin{proposition}[Generality up to constant gauge]\label{prop:general}
Let $A=\sum_iM_i\hat\theta^i$ be a left-invariant $SU(2)$ field on $(t,1,1)$ that is invariant under the right circle
generated by $e_0$, lifted by a constant gauge circle of weight $k$ about a Lie axis. Then $A$ is constant-gauge-equivalent
to a member of the family with $y\ge0$. At $y=0$ the members $(x,0)$ and $(-x,0)$ are gauge-equivalent.
\end{proposition}
\begin{proof}
Write $M=(M_{ia})$ (frame $i$, Lie $a$) and take the gauge axis to be $\tau_0$. The infinitesimal condition is
$LM-kML=0$, with $L$ the $\mathfrak{so}(2)$ generator on the indices $1,2$. Its solution space (the commutant), found
symbolically, has dimension $1$, $3$, $3$, $3$, $1$ for $k=-2,\dots,2$. For $k=\pm2$ it is $m_0=x\,\tau_0$,
$m_1=m_2=0$. For $k=0$ it is $m_0$ arbitrary, $m_1=m_2=0$. For $k=+1$ it is $m_0=x\tau_0$, $m_1=y\tau_1-z\tau_2$,
$m_2=z\tau_1+y\tau_2$, and for $k=-1$ it is $m_0=x\tau_0$, $m_1=-y\tau_1+z\tau_2$, $m_2=z\tau_1+y\tau_2$. For
$|k|\ge3$ and for half-integer $k$ only the weight-$0$ solution with $y=0$ remains. Constant gauge rotations reduce each
case. For $k=+1$, the rotation by $\arg(y+iz)$ about $\tau_0$ maps $(x,y,z)$ to $(x,\sqrt{y^2+z^2},0)$. For $k=-1$,
a rotation by $\pi$ about the first axis maps the $(x,y)$ member to the family member at $(-x,y)$. For $k=0$, a
rotation maps $m_0$ to $|m_0|\tau_0$, a member with $y=0$; a rotation by $\pi$ about $\tau_1$ shows $(x,0)\sim(-x,0)$.
The gauge maps were checked on the fields to $2.2\times10^{-16}$ (\path{scripts/edge_wp4_design_checks.py}, receipt
\path{research/edge-theorem-20261121/wp4_design_checks.json}).
\end{proof}

The proposition says nothing outside its class. It does not cover fields that are circle-invariant but not
left-invariant, lifts by non-constant gauge transformations, or circles other than the right circle about $e_0$. For
$t\ne1$ it assumes, without re-proving it, that the isometry group is $SU(2)_L\times U(1)_R$. As a code control
(computed), full spectra of $H+GG^*$ in the sectors $j\le4$ at six random points agree after removing $z$ to
$1.95\times10^{-15}$, after the $k=-1$ reduction to $x\mapsto-x$ exactly ($0.0$), and after a random constant gauge
frame to $2.23\times10^{-15}$. At $y=0$ the spectrum is therefore even in $x$.

\subsection{Charge blocks}\label{subsec:blocks}

On $V_j$ take the weights $m\in\{-l,\dots,l\}$, $l=j/2$, along the long axis. Write the frame index and the Lie index in
the eigenbasis of the rotation, with charges $c_f,c_l\in\{-1,0,1\}$.

\begin{lemma}[Charge blocks]\label{lem:blocks}
For every member of the family, the charge $q=m+c_f+c_l$ is conserved by $H+GG^*$. Each sector is the orthogonal sum
of charge blocks $B_{j,q}$, spanned by the at most nine states $(m,c_f,c_l)$ with $m=q-c_f-c_l$, $|m|\le l$. An
antiunitary map commutes with $H+GG^*$ and reverses $q$, so the blocks $B_{j,q}$ and $B_{j,-q}$ have the same spectrum.
\end{lemma}
\begin{proof}
$\rho_0$ is diagonal and $\rho_1\pm i\rho_2=-2iJ_\pm$ shift $m$ by $\pm1$; $\Omega_0$ is diagonal and $\Omega_\pm$ shift
the charge $c=c_f+c_l$ by $\pm1$. So $X_0$ preserves $q$ and $X_\pm$ shift it by $\pm1$, and
$\sum_iX_i^*X_i=X_0^*X_0+\tfrac12(X_+^*X_++X_-^*X_-)$ preserves $q$. The curvature term $R_0$ preserves $c$ and does not
touch $m$. For the mirror, take $J\otimes K$, with $J$ the real ($j$ even) or quaternionic ($j$ odd) structure of $V_j$
and $K$ complex conjugation on the frame and Lie indices. This is the argument of~\cite[Lemma ``Charge
blocks'']{davisfloor2026} with $x$ and $y$ free (\path{scripts/ricci_su2_exact_edge.py}; on the Berger path the exact
check is \path{research/su2-floor-family-20261119/charge_selection_symbolic.txt}).
\end{proof}

\paragraph{Closed-form rows.} The rows of every block are symbolic in $(x,y,t)$. For a state $(m,a)$, $a=(c_f,c_l)$,
the diagonal entry is
\begin{equation}\label{eq:diag}
  D(m,a)=\Bigl(w_a-\frac{2m}{\sqrt t}\Bigr)^2+4\bigl(l(l+1)-m^2\bigr)+C_{aa},\qquad C=-(\Omega_1^2+\Omega_2^2)+R_0,
\end{equation}
and the ladder couplings are $2i\sqrt{(l-m)(l+m+1)}\,(\Omega_-)_{a'a}$ to $m+1$ and
$2i\sqrt{(l+m)(l-m+1)}\,(\Omega_+)_{a'a}$ to $m-1$. Here $\Omega_i$ is the action of $M_i$ on the Lie index together
with the Levi-Civita part on the frame index, written in the charge basis. The fibre rotations $w_a$ are
\begin{center}
\begin{tabular}{@{}cccc@{}}
\toprule
$(c_f,c_l)$ & $w_a$ & $(c_f,c_l)$ & $w_a$\\
\midrule
$(1,1)$ & $\sqrt t-x-2/\sqrt t$ & $(0,0)$ & $0$\\
$(1,0)$ & $(t-2)/\sqrt t$ & $(0,-1)$ & $x$\\
$(1,-1)$ & $\sqrt t+x-2/\sqrt t$ & $(-1,1)$ & $-\sqrt t-x+2/\sqrt t$\\
$(0,1)$ & $-x$ & $(-1,0)$ & $(2-t)/\sqrt t$\\
 & & $(-1,-1)$ & $-\sqrt t+x+2/\sqrt t$\\
\bottomrule
\end{tabular}
\end{center}
Read off the table, $w_a=c_f(t-2)/\sqrt t-c_l\,x$. Every entry of $C$ has degree at most one in $x$, and the ladder
couplings $\Omega_\pm$ do not involve $x$ (exact; receipt \path{wp4_design_checks.json}). On the Berger path
$w_a=c\,(t-2)/\sqrt t$ and every nonzero entry of $\Omega_\pm$ has modulus $\sqrt{2t}$
(\path{charge_selection_symbolic.txt}). In double precision the blocks reproduce the full sector spectra at $12$ random
points to $1.47\times10^{-15}$ (\path{research/su2-floor-family-20261119/exact_edge.json}), and the closed-form rows
agree with the general sector builder to $1.7\times10^{-13}$ (receipt \path{family_block_cert_j80.json}).

% ---------------------------------------------------------------------------------------------------
\section{The twist and its closed form}\label{sec:twist}

\begin{proposition}[The twist is exact]\label{prop:exact}
At every member of the family and in every sector $j$, the twist $b=f\,(\hat\theta^1+i\hat\theta^2)\otimes(\tau_1+i\tau_2)$
is an eigenvector of $H+GG^*$ and of $H$, and it is gauge-orthogonal. Its eigenvalue has multiplicity at least two in
the full spectrum, shared with its mirror.
\end{proposition}
\begin{proof}
In charge coordinates $b$ is $(m,c_f,c_l)=(l,1,1)$, of charge $q=l+2$. It is the unique state of maximal charge, so it
spans a $1\times1$ block and is an eigenvector of $H+GG^*$. A gauge parameter $\xi$ carries only a weight and a Lie
index, so the charges in the range of $G$ reach at most $l+1$. Hence $G^*b=0$ and $Hb=(H+GG^*)b$. By
Lemma~\ref{lem:blocks}, the block of charge $-l-2$ has the same entry. This is the counting of~\cite{davisfloor2026}
with $x$ and $y$ free.
\end{proof}

In the computations the multiplicity is exactly two where checked: in the sectors $j\le4$ at
$(x,y,t)=(-2,1.5,4)$ and $(-1,0.7,6)$, in both the full and the gauge-orthogonal spectra. Over $x\in[-40,1]$,
$t\in\{3,5,10\}$, $y\in\{0.5,1.5,3\}$, $j\le12$, the mirror copy is present: the mirror's value minus $\lambda$ is at least
$-6.8\times10^{-13}$ (double precision). A single-copy removal of the twist is therefore an error; every edge computation
removes two copies.

\begin{proposition}[The closed form]\label{prop:closed}
For every $t>0$, $x,y\in\RR$ and $j\ge0$, the twist's eigenvalue is
\begin{equation}\label{eq:closed}
  \lambda(x,y,t,j)=\Bigl(x+\frac{j+2}{\sqrt t}\Bigr)^2+2\sqrt t\,x-y^2 .
\end{equation}
On the Berger field it equals $\lambda_t(j)$ exactly. On $A=s\theta$ it equals $(j+2s+2)^2/t+4s(1-s)$.
\end{proposition}
\begin{proof}
The twist is alone in its charge block, so its eigenvalue is the diagonal entry of $H+GG^*$ there: by
\eqref{eq:weitz}, $\norm{X_0b}^2+\norm{X_1b}^2+\norm{X_2b}^2+\langle\beta,R_0\beta\rangle$ for $\norm b=1$, with $\beta$
its internal frame. Equivalently, take \eqref{eq:diag} at $m=l$, $a=(1,1)$, with $w=\sqrt t-x-2/\sqrt t$ and
$C=4\sqrt t\,x-t-y^2+4$. Then $(w-j/\sqrt t)^2=(x+(j+2)/\sqrt t)^2-2\sqrt t\,x-2(j+2)+t$, the second term of
\eqref{eq:diag} is $4l=2j$, and the sum is \eqref{eq:closed}. Substituting $x=(2-t)/\sqrt t$, $y=\sqrt t$ gives
$\lambda_t(j)$ with difference $0$, and substituting $x=2s/\sqrt t$, $y=2s$ gives the second formula. The computation is
symbolic in $x$, $y$, $t$ (\path{scripts/ricci_su2_twist_family_closed_form.py}, receipt
\path{research/su2-floor-family-20261119/twist_family_closed_form.txt}). As a check (computed), it agrees with the
general builder at $40$ random $(x,y,t,j)$ to $1.4\times10^{-15}$.
\end{proof}

\begin{corollary}[The lift and its discount]\label{cor:lift}
The cost of lifting the twist from the sector $0$ to the sector $j$ is
\[
  \lambda(j)-\lambda(0)=\frac{j^2+4j}{t}+\frac{2xj}{\sqrt t}.
\]
For $j\ge1$ the second term lowers the cost if and only if $x<0$, that is, when the field along the long circle turns
against the twist's own rotation.
\end{corollary}
\begin{proof}
Expand \eqref{eq:closed}: $(x+(j+2)/\sqrt t)^2-(x+2/\sqrt t)^2=2xj/\sqrt t+((j+2)^2-4)/t$.
\end{proof}

On the Berger field $x=(2-t)/\sqrt t$, which changes sign at $t=2$. That is the stretch at which, by the Casimir remark
of~\cite{davisfloor2026}, the coupling of the lift to the stretch vanishes: $\lambda_2(j)-\lambda_2(0)=(j^2+4j)/2$, and
$\lambda_2(0)=0$. For $t>2$ the Berger field lies on the side $x<0$, where the discount helps the twist.

\begin{proposition}[Anatomy]\label{prop:anatomy}
Let $\beta=(\hat\theta^1+i\hat\theta^2)\otimes(\tau_1+i\tau_2)$ be the twist's internal frame. For every $x,y,t$ and $j$:
\begin{enumerate}
  \item $X_0b=i\kappa\,b$ with $\kappa=(t-2)/\sqrt t-x-j/\sqrt t$. The three terms are the Levi-Civita rotation of the
    frame, the gauge field acting on the twist's Lie charge (minimal coupling), and the orbital motion along the fibre.
  \item $(\Omega_1+i\Omega_2)\beta=0$, so with the top weight $X_+b=0$: the twist is a lowest Landau level across the
    whole family.
  \item $R_0\beta=(\mathrm{Ric}_\perp-2B)\beta$ with $\mathrm{Ric}_\perp=4-2t$ and $B$ as in \eqref{eq:B}; the part
    $-2[F,\cdot\,]$ contributes $4\sqrt t\,x-2y^2=-2B$.
  \item $\lambda=\kappa^2+(2j+t+y^2)+(\mathrm{Ric}_\perp-2B)$: fibre momentum squared, the transverse lowest-Landau-level
    term, and the curvature term (Ricci across the fibre, and spin against the fibre field).
  \item $\partial\lambda/\partial x=-2\kappa+4\sqrt t$. The strength $x$ enters twice: through $\kappa$ by minimal
    coupling, and through $B$ by the spin.
  \item On the Berger path the gauge part of $\kappa$ equals the frame part, and $\kappa=(2t-4-j)/\sqrt t$, the transverse
    term is $2(j+t)$, the curvature term is $12-8t$, and $B=3t-4$.
\end{enumerate}
\end{proposition}
\begin{proof}
Exact symbolic computation of each identity; the parts in (4) minus \eqref{eq:closed} is $0$
(\path{scripts/ricci_su2_twist_anatomy.py}, receipt \path{research/su2-floor-family-20261119/twist_anatomy.json}).
Item (5) is the derivative of \eqref{eq:closed}: $-2\kappa+4\sqrt t=2(x+(j+2)/\sqrt t)+2\sqrt t$. Item (6) recovers the
proposition of~\cite{davisfloor2026} that the twist is a lowest Landau level on the Berger field. Item (4) and the
Berger slice of (6) are also proved in Lean~4 (\texttt{price\_anatomy}, \texttt{kappa\_berger},
\texttt{fieldCurv\_berger}, \texttt{price\_berger}; Section~\ref{subsec:lean}).
\end{proof}

Near the edge the twist's slope in $x$ is about $10$ at every point computed ($9.74$, $9.62$, $10.18$ at the three
points of Section~\ref{subsec:competitor}; computed). The competitor's is much smaller (Section~\ref{subsec:competitor}).

\begin{remark}[A reading]\label{rem:landau}
The structure of Proposition~\ref{prop:anatomy} resembles a lowest-Landau-level kinetic term plus a Zeeman term with
gyromagnetic factor $2$. That is the flat-space mechanism of the unstable mode of Nielsen and
Olesen~\cite{nielsenolesen1978}, stated explicitly in~\cite{bcgk2014}. Landau levels on spheres are a standard
construction~\cite{haldane1983,wuyang1976}. I note the resemblance only; no limit relating the two is proved, and none
of these sources says anything about Berger spheres.
\end{remark}

\begin{remark}[The split of $R_0$]\label{rem:split}
The identity $R_0=\mathrm{Ric}\otimes I-2[F,\cdot\,]$ was checked numerically to $4\times10^{-15}$ on the Berger field
and the squashed meron at $t=1,2,3,4,6,8$. On the twist's internal frame the Ricci share is $4-2t$ on both fields.
The field-strength share is $8-6t$ on the Berger field ($-4$ at $t=2$, $-40$ at $t=8$), but positive on the squashed
meron: $+1.2614$ at $t=2$, $+0.9265$ at $t=3$, $+0.7334$ at $t=4$, $+0.5183$ at $t=6$, $+0.4010$ at $t=8$ (computed;
\path{research/su2-floor-family-20261119/gyromagnetic_split.txt}). Read through item (3) of
Proposition~\ref{prop:anatomy}, on the squashed meron the field strength is opposed to the twist's spin.
Section~\ref{sec:decide} shows that this alone does not decide the floor.
\end{remark}

% ---------------------------------------------------------------------------------------------------
\section{The edge}\label{sec:edge}

\subsection{Edge functions}\label{subsec:edgefunctions}

By Lemma~\ref{lem:blocks}, whether the twist is the floor of the sectors $j\le J$ is a comparison between $\lambda$ and
the lowest eigenvalues of finitely many blocks. For $H+GG^*$ define
\[
  \phi_A(x)=\min_{j\le J}\ \min_{q\ne\pm(l+2)}\bigl[\lambda_{\min}(B_{j,q})-\lambda_j\bigr],
\]
with both twist blocks excluded, and for the gauge-orthogonal spectrum $\phi_S(x)=\min_{j\le J}[\mu_{\min,j}-\lambda_j]$,
where $\mu$ is the slice spectrum from the general builder after removing two copies of $\lambda_j$. The edge $x^*_A$
(resp.\ $x^*_S$) is the largest $x$ in the scan window where $\phi$ changes sign from $\ge0$ to $<0$, located by Brent's
method to $10^{-12}$; $x^*_{12}$ denotes the $H+GG^*$ edge with $J=12$. At the edge $\lambda_j$ is an eigenvalue of the
competitor block, so the edge is a root of $\det(B_{j,q}(x,y,t)-\lambda_j)=0$. The gauge-orthogonal edge is slightly
higher than the $H+GG^*$ edge at every point of Table~\ref{tab:edge}: at $t=3$ by $0.027$ ($y=\sqrt3$) and $0.020$
($y=3/2$). In double precision the gauge-orthogonal edges agree with the independent bisection of the earlier edge test
at its $21$ points to $3\times10^{-5}$, within that bisection's tolerance $10^{-4}$ (\path{scripts/ricci_su2_exact_edge.py},
receipt \path{exact_edge.json}).

\begin{table}[t]
\centering
\caption{The edge in the sectors $j\le12$, $x^*_A$ / $x^*_S$ (double precision; not certified). For $y\ge1.5$ and
$y=\sqrt t$ the competitor is $B_{4,0}$ throughout. Every positive entry is an artefact of the cut-off $j\le12$
(Section~\ref{sec:interval}). Receipt \texttt{research/su2-floor-family-20261119/}\allowbreak\texttt{exact\_edge.json}.}
\label{tab:edge}
\begin{tabular}{@{}cccc@{}}
\toprule
$t$ & $y=0.5$ & $y=1.5$ & $y=\sqrt t$\\
\midrule
$2$ & $-0.613240$ / $-0.605292$ & $-0.682793$ / $-0.657567$ & $-0.677814$ / $-0.655100$\\
$3$ & $-0.335488$ / $-0.318061$ & $-0.371582$ / $-0.351580$ & $-0.382168$ / $-0.355187$\\
$4$ & $-0.202972$ / $-0.191702$ & $-0.208048$ / $-0.191652$ & $-0.221163$ / $-0.189150$\\
$6$ & $-0.082842$ / $-0.078032$ & $-0.044874$ / $-0.033249$ & $-0.035386$ / $+0.008031$\\
$8$ & $-0.033799$ / $-0.031530$ & $+0.031482$ / $+0.040169$ & $+0.082513$ / $+0.138269$\\
\bottomrule
\end{tabular}
\end{table}

Table~\ref{tab:edge} gives the edge at five stretches. At $t=2$ and $y=0.5$ the competitor is $B_{3,1/2}$; at $t=3$,
$y=0.5$ it is $B_{2,-1}$, the mirror of $B_{2,1}$. Within $j\le12$ the edge rises with the stretch and turns positive for
strongly stretched spheres with a strong cross field. That is an artefact of the cut-off: Section~\ref{sec:interval}
proves that for $x>0$ high sectors undercut the twist.

\subsection{The competitor}\label{subsec:competitor}

For a cross field $y$ around $1$ and above, the edge is set by $B_{4,0}$, the $9\times9$ charge-$0$ block of the sector
$j=4$, with entries in $\mathbb Q(\sqrt2,\sqrt3)$. (In the map of Section~\ref{subsec:competitors}, for $t\le6$ its
region runs from $y\approx0.75$--$0.99$ up to a curve from $y\approx1.9$ at $t=2$ to $y\approx3$ at $t\approx5.7$.) At
the three points computed, its lowest state is a mirror-symmetric superposition dominated by
$(m,c_f,c_l)=(+2,-1,-1)$ and $(-2,+1,+1)$: orbit and internal spin opposed, total fibre charge zero. The twist is
charged, with orbit and spin aligned. The competitor is nearly blind to $x$. In double precision:
\begin{center}
\small
\begin{tabular}{@{}lccc@{}}
\toprule
point & weight in $m=\pm2$ & slope of competitor & slope of twist ($j=4$)\\
\midrule
Berger at $t^*$ ($x=-0.42496$, $y=1.64257$, $t=2.69803$) & $0.62$ & $-1.561$ & $9.741$\\
$t=4$, $y=2$, $x=-0.19$ & $0.66$ & $-0.623$ & $9.620$\\
$t=8$, $y=\sqrt8$, $x=0.14$ & $0.73$ & $+1.684$ & $10.179$\\
\bottomrule
\end{tabular}
\end{center}
The function $\lambda_{\min}(B_{4,0})-\lambda(x,y,t,4)$ is strictly decreasing on $x\in[-3,3]$ at all $48$ pairs
$(t,y)$ tested, $t$ from $2$ to $10$ and $y$ from $0.5$ to $3$ (computed; receipt \path{twist_anatomy.json}). An
operator-norm bound on $\partial_xB_{4,0}$ is not enough to prove this: it works at $21$ of $27$ edge points and fails at
every point with $y=\sqrt t$, $t\ge4$.

\subsection{The edge polynomial}\label{subsec:edgepoly}

\begin{proposition}[The edge polynomial]\label{prop:edgepoly}
Let $s=\sqrt t$ and let $P(x,y,s)$ be the numerator of $\det\bigl(B_{4,0}(x,y,s^2)-\lambda(x,y,s^2,4)\,I\bigr)$. Then $P$
has rational coefficients, $473$ terms, degrees $(12,14,28)$ in $(x,y,s)$ and total degree $32$, and
\[
  P=-16384\,P_1P_2 ,
\]
with $P_1$, $P_2$ irreducible over $\mathbb Q$, each of multiplicity one and each even in $y$. $P_1$ has degrees
$(5,6,12)$ in $(x,y,s)$ and $61$ terms; $P_2$ has degrees $(7,8,16)$ and $113$ terms. They begin
\begin{align*}
  P_2&=-8s^{16}x^2+24s^{15}x^3+8s^{15}xy^2-4s^{14}x^4-42s^{14}x^2y^2-320s^{14}x^2-2s^{14}y^4-11s^{13}x^5+\cdots,\\
  P_1&=2s^{12}x^2-5s^{11}x^3-3s^{11}xy^2-24s^{11}x-s^{10}x^4+\cdots .
\end{align*}
\end{proposition}
\begin{proof}
Exact symbolic determinant over $\mathbb Q(\sqrt2+\sqrt3)[x,y,s]$ and exact factorization~\cite{sympy2017}
(\path{scripts/ricci_su2_edge_polynomial.py}; full polynomials in
\path{research/su2-floor-family-20261119/edge_polynomial.txt} and \path{edge_polynomial.srepr}). The factorization is
also proved in Lean (Theorem~\texttt{edgeDet\_factor}, Section~\ref{subsec:lean}). A second, independent builder of
the blocks recovers the same edge factor exactly (quotient $1$, remainder $0$;
\path{research/edge-theorem-20261121/wp1_block_polynomials.txt}). Its full determinant was compared with
\path{edge_polynomial.srepr} only by the ratio $1$ at four rational points, which is a computed check, not a proof.
\end{proof}

$P_2$ vanishes on the numerically located edge wherever $B_{4,0}$ is the competitor: at all $77$ grid edges of the map
the nearest root of $P_2$ lies within $2.4\times10^{-13}$ of the edge, while the scale-free residual of $P_1$ there is
at least $2.7\times10^{-3}$, and at ten sample points the nearest root of $P_1$ is $0.36$--$0.55$ away (computed).
\emph{The edge is not in general the largest real root of $P_2$.} On the $18$ known-answer rows of
Table~\ref{tab:edge}'s receipt with $y=1.5$ or $y=\sqrt t$, $P_2$ has three or four further real roots in $[-3,3]$ above
the edge, from about $0.55$ ($t=2$) to $3.5$ ($t=12$, $y=\sqrt t$), where a higher eigenvalue of $B_{4,0}$ meets
$\lambda_4$. The edge is the root at which the \emph{lowest} eigenvalue of $B_{4,0}$ equals $\lambda_4$. This was found
by a known-answer gate of a pre-registered test, before any of its held-out points was computed
(Section~\ref{sec:prereg}); the earlier statement ``largest real root'' is withdrawn.

\subsection{The Berger crossings}\label{subsec:berger}

\begin{proposition}[Berger restriction]\label{prop:bergerrestriction}
For $s\ne0$,
\[
  P_2\Bigl(\frac{2-s^2}{s},s,s\Bigr)=-32\,(s-1)(s+1)\,Q_{16}(s),
\]
\[
  Q_{16}(s)=3s^{16}+121s^{14}+1066s^{12}-2836s^{10}-18624s^8+20992s^6+181248s^4-425984s^2+262144 ,
\]
the polynomial of~\cite[Remark ``Below the range'']{davisfloor2026} whose largest positive root is $s=\sqrt{t_A}$,
$t_A=2.7311222782381135$. So $t_A$ is algebraically a root of the edge polynomial on the Berger path.
\end{proposition}
\begin{proof}
Exact polynomial division with zero remainder (\path{scripts/ricci_su2_edge_polynomial.py}, receipt
\path{edge_polynomial.txt}). Also proved in Lean (\texttt{edgeFactor\_berger}), together with a root of $Q_{16}$ in
$[1.65261,1.65262]$ and a zero of the restriction with $s^2\in[2.73111,2.73116]$.
\end{proof}

Along the Berger path $x=(2-t)/\sqrt t$, $y=\sqrt t$, the computed edges are crossed as follows (double precision,
root-finding to $10^{-12}$ on exact ingredients; receipt \path{exact_edge.json}). The $H+GG^*$ edge is crossed at
$t=2.7311222782381144$, which is $t_A$ to $9\times10^{-16}$. The gauge-orthogonal edge is crossed at
$t=2.6980269405211303$, which is $t^*=2.6980269405200263$ to $1.1\times10^{-12}$. Both crossings go through the sector
$j=4$, the first through the block $B_{4,0}$. That $t_A$ is a root of $P_2$ on the path is proved
(Proposition~\ref{prop:bergerrestriction}); that the path crosses the edge there is computed. Between $t^*$ and $t_A$ the
eigenvalue of $H+GG^*$ below $\lambda_t(4)$ is not gauge-orthogonal (computed; \cite[Remark ``Below the
range'']{davisfloor2026}).

The thresholds of~\cite{davisfloor2026} are therefore where the Berger path meets the edge, not properties of the Berger
field alone. The edge is a crossing, not an avoided one: the twist is alone in its charge block, every competitor found
has a different charge, and eigenvalues of different symmetry type may cross where those of one type generically do
not~\cite{vonneumannwigner1929}.

\subsection{Other competitors}\label{subsec:competitors}

On a grid of $132$ points, $t=2,\dots,12$ by $y=0.25,\dots,3$, with $j\le12$ and the window $[-3,1]$, eight blocks (up
to the mirror $q\leftrightarrow-q$) set the $H+GG^*$ edge (computed). Their edge factors were derived exactly
(Table~\ref{tab:competitors}); the determinants and factorizations over $\mathbb Q$ are exact algebra, and the edge
factors of $B_{1,3/2}$, $B_{2,2}$, $B_{12,7}$ and $B_{2,1}$ were checked irreducible over $\mathbb Q$ again by an
independent rerun. All $132$ grid edges are roots of the right factor, to $2.4\times10^{-13}$ (computed). At every grid
point $\phi_A$ has a single sign change on a $401$-point scan of $[-3,1]$, and $\phi_A(-3)\ge1.43$.

\begin{table}[t]
\centering
\small
\caption{Blocks that set the edge in the sectors $j\le12$ (map computed; polynomials exact). Degrees are in
$(x,y,s)$; ``$s\cdot$'' marks a factor $s$ split off. $\dagger$: block of the top sector allowed, a cut-off
artefact. Receipts \texttt{research/edge-theorem-20261121/}\allowbreak\texttt{wp1\_competitor\_map.md} and
\texttt{wp1\_block\_polynomials.txt}.}
\label{tab:competitors}
\begin{tabular}{@{}lclcccc@{}}
\toprule
block & size & entries in & edge factor & terms & grid edges & max $|\text{root}-\text{edge}|$\\
\midrule
$(4,0)$ & $9$ & $\mathbb Q(\sqrt2,\sqrt3)$ & $P_2$, $(7,8,16)$ & $113$ & $77$ & $2.4\times10^{-13}$\\
$(3,\tfrac12)$ & $8$ & $\mathbb Q(\sqrt2,\sqrt6)$ & $(11,14,26)$ & $400$ & $6$ & $1.4\times10^{-14}$\\
$(2,1)$ & $6$ & $\mathbb Q$ & $s\cdot(8,10,18)$ & $169$ & $20$ & $2.3\times10^{-13}$\\
$(1,\tfrac12)$ & $5$ & $\mathbb Q(\sqrt2)$ & $s\cdot(7,10,16)$ & $133$ & $4$ & $1.5\times10^{-16}$\\
$(1,\tfrac32)$ & $3$ & $\mathbb Q(\sqrt2)$ & $s\cdot(4,4,8)$ & $26$ & $15$ & $8.5\times10^{-14}$\\
$(2,2)$ & $3$ & $\mathbb Q$ & $s\cdot(4,4,8)$ & $27$ & $2$ & $1.7\times10^{-14}$\\
$(12,6)^\dagger$ & $6$ & $\mathbb Q(\sqrt6,\sqrt{11})$ & $s\cdot(8,10,18)$ & $172$ & $1$ & $1.7\times10^{-15}$\\
$(12,7)^\dagger$ & $3$ & $\mathbb Q(\sqrt6)$ & $s\cdot(4,4,8)$ & $27$ & $7$ & $1.1\times10^{-13}$\\
\bottomrule
\end{tabular}
\end{table}

The regimes are ordered by the cross field. At small $y$ the competitor is $B_{2,1}$, up to $y\approx0.57$--$0.67$;
then $B_{3,1/2}$, up to $y\approx0.75$--$0.99$; then $B_{4,0}$; and at large $y$ and small $t$, $B_{1,3/2}$ (and $B_{2,2}$
at $t=2$, $y\ge2.75$). For $t\ge7$ at $y\approx0.75$--$1$ the blocks $B_{1,1/2}$ and $B_{1,3/2}$ compete. Examples of
switches (computed): from $B_{3,1/2}$ to $B_{2,1}$ along $y=0.5$ at $t=2.304727$ (edge $-0.49703$); from $B_{3,1/2}$ to
$B_{4,0}$ at $t=2$, $y=0.752443$ (edge $-0.611741$); from $B_{4,0}$ to $B_{1,3/2}$ at $t=2$, $y=1.896879$ (edge
$-0.699953$). At $t=3,4,5,6$ the switch from $B_{2,1}$ to $B_{3,1/2}$ is at $y=0.574998$, $0.621151$, $0.650418$,
$0.673813$, and from $B_{3,1/2}$ to $B_{4,0}$ at $y=0.874254$, $0.930635$, $0.965022$, $0.988971$. The closest runner-up
gap is $0.0053$, at $(t,y)=(2,0.75)$. One boundary, from $B_{12,6}$ to $B_{2,1}$ along $y=0.25$ for $t\in[2,3]$, was not
located. The smallest of these blocks has a short edge factor; for $B_{1,3/2}$ it is
\begin{multline*}
2s^8x-5s^7x^2-s^7y^2-s^6x^3+7s^6xy^2+s^5x^4+2s^5x^2y^2-23s^5x^2-s^5y^4+7s^5y^2-3s^4x^3y^2\\
+9s^4x^3+3s^4xy^4-7s^4xy^2-30s^4x+2s^3x^4-10s^3x^2y^2+4s^3x^2+6s^3y^4+2s^3y^2\\
+20s^2x^3-28s^2xy^2-20s^2x+56sx^2-40sy^2+48x .
\end{multline*}
Six of the eight blocks belong to two near-top families of charge $q=l+1$ and $q=l$; only $B_{4,0}$ and $B_{3,1/2}$
are sporadic. Those families govern the high sectors, and Section~\ref{sec:interval} uses them.

\subsection{Lean formalisation}\label{subsec:lean}

The algebraic identities above, and the closed form's anatomy, are proved in Lean~4~\cite{lean4} (v4.36.0) with
Mathlib~\cite{mathlib2020}: sixteen theorems, with no \texttt{sorry} and no axioms beyond \texttt{propext},
\texttt{Classical.choice} and \texttt{Quot.sound} (\path{docs/papers/berger-spectrum/lean/Helike/Edge.lean} and
\path{Price.lean}).

\paragraph{\texttt{Helike/Edge.lean} (four theorems).} The polynomials are those of the receipt
\path{edge_polynomial.srepr}; the $473$-term determinant numerator is entered in ten parts of at most $48$ terms, checked
equal to the receipt by sympy before the Lean was written.
\begin{itemize}
  \item \texttt{edgeDet\_factor}: for all real $x,y,s$, $\mathtt{edgeDet}=-16384\,P_1P_2$.
  \item \texttt{edgeFactor\_berger}: for $s\ne0$, $P_2((2-s^2)/s,s,s)=-32(s-1)(s+1)Q_{16}(s)$.
  \item \texttt{tA\_poly\_root}: $Q_{16}$ has a root $s\in[1.65261,1.65262]$, by the intermediate value theorem
    ($Q_{16}(1.65261)<0<Q_{16}(1.65262)$).
  \item \texttt{berger\_crosses\_edge}: there is $s\in[1.65261,1.65262]$ with $P_2((2-s^2)/s,s,s)=0$ and
    $s^2\in[2.73111,2.73116]$.
\end{itemize}

\paragraph{\texttt{Helike/Price.lean} (twelve theorems).} With $t=s^2$, $s\ne0$, the closed form \eqref{eq:closed} is
\texttt{price}, $\kappa$ is \texttt{kappa}, $B$ is \texttt{fieldCurv}, and $\lambda_t(j)$ is
$\mathtt{floorLaw}\,t\,j=(j^2+(8-2t)j-2t^2-4t+16)/t$. The theorems are \texttt{price\_anatomy} (item (4) of
Proposition~\ref{prop:anatomy}); \texttt{floorLaw\_paper\_form}; \texttt{kappa\_berger}, \texttt{fieldCurv\_berger} and
\texttt{price\_berger} (the Berger slice, and $\lambda=\lambda_t(j)$ there); \texttt{floorLaw\_three},
\texttt{floorLaw\_three\_sector0} ($-14/3$) and \texttt{floorLaw\_three\_sector12} ($154/3$) at $t=3$;
\texttt{floorLaw\_vertex\_form}, $\lambda_t(j)=(j-(t-4))^2/t+(4-3t)$, with \texttt{floorLaw\_ge\_vertex} and
\texttt{vertex\_on\_line}; and \texttt{floorLaw\_cutoff}, $\lambda_t(t-4+\sqrt{t(3t-4)})=0$ for $t\ge4/3$.

The builds completed on two machines (Price: $18254$ jobs; a full rebuild: $18256$ jobs, with \texttt{Helike.Edge}
taking $231$\,s for the factorization, of degree $28$ in $s$).

\paragraph{What Lean does not check.} Lean checks polynomial identities and the two brackets. It does not check that
\texttt{edgeDet} is the determinant of $B_{4,0}-\lambda_4I$ (a computer-algebra step), that the root of $P_2$ is the
edge (an eigenvalue comparison), that the twist is the floor anywhere, or that \texttt{price} is an eigenvalue of the
Hessian. Those rest on Propositions~\ref{prop:exact}--\ref{prop:edgepoly} and on the certificates of
Sections~\ref{sec:interval} and~\ref{sec:domain}.

% ===================================================================================================
\section{The region is an interval, and $x>0$ loses}\label{sec:interval}

Section~\ref{sec:edge} locates the edge in the sectors $j\le12$. Two facts change the picture once all sectors and
all $x$ are allowed. On the side of positive $x$ the twist always loses, to high sectors. On the side of very negative
$x$ it loses too, to a state of the sector $j=0$. Its region in $x$ is therefore bounded on both sides.

Throughout, $l=j/2$, and a row of a charge block is labelled by its weight $m=l-\xi$ and its type $a=(c_f,c_l)$, as in
Section~\ref{sec:setting}. Write $\lambda_j=\lambda(x,y,t,j)$ for the twist's eigenvalue of
Proposition~\ref{prop:closed}.

\subsection{Positive $x$}

\begin{theorem}[For $x>0$ the twist is not the floor of every sector]\label{thm:xpos}
Let $t>0$, $x,y\in\mathbb{R}$ and $l\in\frac12\mathbb{Z}_{\ge0}$.
\begin{enumerate}
  \item In the charge block $B_{2l,l}$, the state $e_-$ of weight $m=l$ and type $(1,-1)$ satisfies
  \[
    \langle e_-,(H+GG^*)e_-\rangle-\lambda_{2l}
      =-\frac{8lx}{\sqrt t}-4\sqrt t\,x+4y^2-\frac{8x}{\sqrt t}.
  \]
  \item In the charge block $B_{2l,l+1}$, the state $e_0$ of weight $m=l$ and type $(1,0)$ satisfies
  \[
    \langle e_0,(H+GG^*)e_0\rangle-\lambda_{2l}
      =-\frac{4lx}{\sqrt t}-2\sqrt t\,x-x^2+3y^2-\frac{4x}{\sqrt t}.
  \]
  \item Consequently, if $x>0$, the twist is not the lowest eigenvalue of $H+GG^*$ in the sector $j=2l$ for every
  \[
    l>\frac{\sqrt t\,\bigl(4y^2-4\sqrt t\,x-8x/\sqrt t\bigr)}{8x}
  \]
  (for every $l$ if the right side is negative). In particular, for $x>0$ the twist is not the floor of every sector.
\end{enumerate}
\end{theorem}

\begin{proof}
The diagonal entries of every charge block are known in closed form (Section~\ref{sec:setting}). For the row
$m=l-\xi$ of type $a$ they satisfy the exact identity
\[
  D_{m,a}-\lambda_{2l}=4\xi(2l-\xi)\Bigl(1-\frac1t\Bigr)+l\,K_a+E_a(\xi)+R_a,\qquad
  K_a=4-\frac8t-\frac{4(w_a+x)}{\sqrt t},
\]
in which the remaining terms $E_a(\xi)$ and $R_a$ do not involve $l$. For the twist's own type, $K_{(1,1)}=0$. The two
types next to it have $K_{(1,0)}=-4x/\sqrt t$ and $K_{(1,-1)}=-8x/\sqrt t$, and at $\xi=0$ the identity gives the two
displayed formulas. Both states are unit vectors of charge $l$ and $l+1$, while the twist has charge $l+2$, so they
lie in blocks other than the twist's. A one-vector Rayleigh quotient bounds the lowest eigenvalue of a block above by
any of its diagonal entries. For $x>0$ the coefficient of $l$ in item~1 is $-8x/\sqrt t<0$, which gives item~3.
The identity is verified symbolically for all nine types, and the Rayleigh bound for both states, with $l$, $x$, $y$
and $t$ symbolic (\path{scripts/edge_wp2_symbolic.py}, receipt
\path{research/edge-theorem-20261121/wp2_symbolic.txt}, ``ALL SYMBOLIC CHECKS PASS''). As a numerical cross-check, the
diagonal of the double-precision block builder minus $\lambda$ matches the formula at $(x,y,t)=(0.1,\sqrt{10},10)$ for
$j=300$, $306$ and $310$ to about $10^{-12}$ (\path{research/edge-theorem-20261121/workflow_reports.json}).
\end{proof}

The theorem is stated for $H+GG^*$. The states $e_-$ and $e_0$ are not shown to be gauge-orthogonal, so for the
gauge-orthogonal spectrum the same conclusion is supported by the computations below but not proved.

\begin{corollary}\label{cor:edgebound}
Let $x^*_{12}(t,y)$ be the upper end of the set of $x$ at which the twist is the floor of $H+GG^*$ in every sector
$j\le12$ (Section~\ref{sec:edge}; it is located numerically on grids). The set of $x$ at which the twist is the floor of
$H+GG^*$ in every sector lies in $x\le\min(x^*_{12},0)$. In short, the all-sector edge is at most
$\min(x^*_{12},0)$.
\end{corollary}

\begin{proof}
The bound $\le0$ is Theorem~\ref{thm:xpos}. The bound $\le x^*_{12}$ holds because the all-sector condition contains
the condition in $j\le12$.
\end{proof}

\paragraph{A worked example (double precision, except the threshold).} At $(x,y,t)=(0.1,\sqrt{10},10)$ the type
$(1,-1)$ has $K=-0.252982212813$ and constant term $38.4821067231$, so by Theorem~\ref{thm:xpos} the twist loses in
every sector with $l>152.113883008$. The type $(1,0)$ has $K=-0.126491106407$ and constant term $29.2310533616$, and
loses for $l>231.091767571$. The actual undercut comes earlier, because the lowest eigenvalue of a block is below its
diagonal. From the blocks, the margin $\min_q\lambda_{\min}(B_{j,q})-\lambda_j$ is $-0.58346$ at $j=120$ (minimising
charge $q=-61$), $-8.9498$ at $j=200$ and $-34.100$ at $j=400$. An independent full-sector builder gives $-0.583$
for $H+GG^*$ and $-0.581$ for the gauge-orthogonal spectrum at $j=120$, and $-8.95$ at $j=200$
(\path{research/edge-theorem-20261121/wp2_large_sector_bound.md}, Section~0, and the verifier's rerun in
\path{workflow_reports.json}).

\paragraph{Near $x=0$ (double precision; not certified).} At $(t,y)=(10,\sqrt{10})$ and $x=0$ the margin stays
positive and rises slowly: $7.7516$ at $j=2000$ and $7.7752$ at $j=20000$. At $x=+0.01$ it is negative at large $j$:
$-7.67$ at $j=2000$ and $-32.9$ at $j=4000$ (\path{workflow_reports.json}). The point $x=0$ is degenerate. There
$K_{(1,0)}=K_{(1,-1)}=0$, the margin at the near end of each sector is of order one, and a proof would need a
second-order argument. No proof below covers $x\in(-0.02,0)$.

\paragraph{The positive edges of the cut-off (double precision; not certified).} Within $j\le12$ the edge turns
positive for strongly stretched spheres with a strong cross field (Section~\ref{sec:edge}). That is an artefact of the
cut-off. At all $58$ grid points of the competitor map with $x^*_{12}>0$, some sector $j>12$ undercuts the twist at
$x^*_{12}$: at $54$ points within $13\le j\le200$ (\path{wp1_tail_check.json}, near-top charges only), and, in the
verifier's rerun, at $(7,2.0)$, $(7,2.25)$ and $(7,2.5)$ by the sectors $j=2000$, $400$ and $400$, and at $(12,3)$ by
$j=100$, with margin $-10.46$. In a full charge scan at four of these points the minimising charge was $|q|=l+1$ or
$l$. Table~\ref{tab:cutoff} shows how the edge falls as the cut-off $J$ is raised at one point. In the other direction,
at $x=-0.01$ and $x=-0.001$, at each of the $57$ positive-edge points of the tail run (all but $(12,3)$), neither the
near-top blocks of $13\le j\le2000$ nor any block of $j\le12$ undercuts the twist; the smallest margin is $0.029$
(\path{research/edge-theorem-20261121/wp1_tail_check.json}, \path{wp1_tail_negx.json}, \path{wp1_fullsector_check.json},
\path{workflow_reports.json}).

\begin{table}[t]\centering\small
\caption{The edge of the sectors $j\le J$ at $(t,y)=(12,2.75)$ (double precision; \texttt{wp1\_competitor\_map.md},
Section~4).}\label{tab:cutoff}
\begin{tabular}{lccccc}\toprule
$J$ & $12$ & $24$ & $48$ & $96$ & $192$\\ \midrule
edge & $0.1825$ & $0.1825$ & $0.1706$ & $0.0969$ & $0.0519$\\
\bottomrule\end{tabular}
\end{table}

\subsection{Very negative $x$}

\begin{proposition}[The lower edge]\label{prop:lower}
Let $v=\hat\theta^0\otimes\tau_0$ in the sector $j=0$ (weight $m=0$, type $(0,0)$, charge $0$). For every
$(x,y,t)$ with $t>0$:
\begin{enumerate}
  \item $G^*v=-[M_0,\tau_0]=0$, so $v$ is gauge-orthogonal; it has charge $0$, so it is orthogonal to both copies of the
  twist (charge $\pm2$ in the sector $0$);
  \item $\langle v,(H+GG^*)v\rangle=\langle v,Hv\rangle=4t+2y^2$, independent of $x$;
  \item the twist is neither the lowest eigenvalue of $H+GG^*$ nor the lowest gauge-orthogonal eigenvalue in the sector
  $j=0$ whenever $x<r_{\mathrm{lo}}$ or $x>r_{\mathrm{hi}}$, where
  \[
    r_{\mathrm{lo},\mathrm{hi}}(t,y)=-\Bigl(\sqrt t+\frac2{\sqrt t}\Bigr)\mp\sqrt{5t+3y^2+4}
  \]
  are the roots of $\lambda_0(x)=4t+2y^2$.
\end{enumerate}
\end{proposition}

\begin{proof}
Item~1: $M_0=x\tau_0$ commutes with $\tau_0$, and $v$ carries no weight, so $G^*v=0$; this holds symbolically for all
$(x,y,t)$ (\path{scripts/edge_wp4_design_checks.py}, receipt \path{wp4_design_checks.json}, entry
\path{Gstar_theta0_tau0_sector0}). Item~2: the row of type $(0,0)$ at $l=0$ in
\path{wp2_symbolic.txt} gives the diagonal entry minus $\lambda_0$ as $-2\sqrt t\,x+4t-x^2+3y^2-4/t-4x/\sqrt t$. Adding
$\lambda_0=(x+2/\sqrt t)^2+2\sqrt t\,x-y^2$ leaves $4t+2y^2$ (exact symbolic check). A rerun measured the constant
value $18.88$ at $(y,t)=(1.2,4)$, with zero second difference in $x$. Item~3: $\lambda_0(x)$ grows like $x^2$, while the
Rayleigh quotient of $v$ does not move. Where $\lambda_0>4t+2y^2$, the vector $v$, orthogonal to the twist and its
mirror and inside the gauge slice, has Rayleigh quotient below $\lambda_0$ for both $H+GG^*$ and the operator on the
slice. Solving the quadratic $\lambda_0(x)=4t+2y^2$ gives $r_{\mathrm{lo},\mathrm{hi}}$.
\end{proof}

Examples: $r_{\mathrm{lo}}(3,0.5)=-7.3308$, $r_{\mathrm{lo}}(3,3)=-9.6691$, $r_{\mathrm{lo}}(10,0.5)=-11.1941$ and
$r_{\mathrm{lo}}(10,3)=-12.7947$. Since $r_{\mathrm{lo}}$ decreases in $y$ and, for $t>2$, in $t$, it follows that
$r_{\mathrm{lo}}\le r_{\mathrm{lo}}(3,0.5)<-7.33$ on the whole domain $D$ of Section~\ref{sec:domain}.

\paragraph{Where the lower edge is (double precision; not certified).} The function $\phi_A$ of
Section~\ref{sec:edge}, scanned on $[-40,-1.5]$, changes sign once. It first turns non-negative at about $x=-4.1$ for
$(y,t)=(3,3)$, $-5.6$ for $(0.5,3)$, $-6.3$ for $(3,10)$, $-7.5$ for $(0.5,10)$ and $-5.75$ for $(1.5,6)$. Every earlier
scan stopped at $x=-3$, which hid the lower end; the region had been read as a half-line.

\begin{corollary}[The region is bounded]\label{cor:interval}
For each $(t,y)$, every $x$ at which the twist is the floor of every sector, for $H+GG^*$ or for the gauge-orthogonal
spectrum, satisfies $r_{\mathrm{lo}}(t,y)\le x\le r_{\mathrm{hi}}(t,y)$; for $H+GG^*$ it also satisfies
$x\le\min(x^*_{12},0)$. Writing $x_{\mathrm{lo}}$ for the lower end of the region, $x_{\mathrm{lo}}\ge r_{\mathrm{lo}}$.
\end{corollary}

\begin{proof}
Proposition~\ref{prop:lower} and Corollary~\ref{cor:edgebound}.
\end{proof}

This is the two-sided condition the title names. That the region is exactly one interval $[x_{\mathrm{lo}},x^*]$, with
no holes, is computed on grids and not proved. What is proved is narrower: Theorem~\ref{thm:domain} shows that the
region contains $[-3,-1.5]$ for every $(t,y)\in[3,10]\times[0.5,3]$, and Theorem~\ref{thm:slice} gives, in the sectors
$j\le12$ on two slices, the exact upper end on $[-3,r_{\mathrm{next}})$.

\subsection{The upper edge on two slices}

\begin{theorem}[The edge on two slices, sectors $j\le12$]\label{thm:slice}
Let $t=3$ and $y\in\{\sqrt3,\,3/2\}$. Assume that the closed-form charge blocks of Section~\ref{sec:setting} are the
charge blocks of $H+GG^*$. Let $x^*$ be the unique real root of $P_2(x,y,\sqrt3)$ in $[-3,1]$, and let $r_{\mathrm{next}}$
be the next real root above it of $\det(B_{4,0}-\lambda_4)$, a root of $P_1(x,y,\sqrt3)$ (Proposition~\ref{prop:edgepoly}).
Then:
\begin{enumerate}
  \item $m(x)=\lambda_{\min}(B_{4,0})-\lambda_4$ is positive on $[-3,x^*)$, zero at $x^*$ and negative on
  $(x^*,r_{\mathrm{next}})$, and $x^*$ is a simple zero;
  \item every charge block $B_{j,q}$ with $j\le12$ other than $B_{4,0}$ and the two twist blocks has
  $\lambda_{\min}(B_{j,q})>\lambda_j$ on $[-3,x^*+0.02]$;
  \item hence, for $x\in[-3,r_{\mathrm{next}})$, the twist is the lowest eigenvalue of $H+GG^*$ in every sector $j\le12$
  if and only if $x\le x^*$. It is the strict floor for $x<x^*$; at $x^*$ the value $\lambda_4$ ties one eigenvalue of
  $B_{4,0}$.
\end{enumerate}
\end{theorem}

\begin{table}[t]\centering\footnotesize
\caption{The two slices of Theorem~\ref{thm:slice} (\texttt{wp3\_slice\_iff.json}). $x^*_A$ is the double-precision
edge of \texttt{exact\_edge.json}.}\label{tab:slice}
\begin{tabular}{lcc}\toprule
 & $y=\sqrt3$ & $y=3/2$\\ \midrule
$x^*$ & $-0.3821684822419167$ & $-0.37158172515020627$\\
$r_{\mathrm{next}}$ & $-0.0650646510788456$ & $-0.008268340671784679$\\
$\partial_xP_2$ on $x^*\pm10^{-6}$ & $\approx1.1032\times10^{7}$ & $\approx8.3613\times10^{6}$\\
$P_1$ on $x^*\pm10^{-6}$ & $[23434.09,\,23434.34]$ & $[20392.46,\,20392.68]$\\
smallest pivot bound, block $(1,3/2)$ & $0.0038$ & $0.027$\\
$x^*-x^*_A$ (computed) & $9.4\times10^{-16}$ & $2.4\times10^{-13}$\\
\bottomrule\end{tabular}
\end{table}

\begin{proof}
\emph{(E1)} The blocks are built exactly in \textsc{SymPy}~\cite{sympy2017}. $B_{4,0}$ is exactly real symmetric, and
the two $1\times1$ blocks of charge $\pm(l+2)$ equal $\lambda_j$ exactly and are set aside. The exact blocks agree with
the double-precision blocks to $2.8\times10^{-14}$.
\emph{(E2)} On the slice, $3^9D=P(x,y,\sqrt3)$ with $D$ the determinant, and the factorization $P=-16384\,P_1P_2$ is
recomputed, with degrees $(5,6,12)$ and $(7,8,16)$, both factors even in $y$.
\emph{(E3)} Write $P_i=A+\sqrt3B$ with $A$, $B$ polynomials in $x$ with rational coefficients. The norms $N=A^2-3B^2$ are squarefree, so every real
root is simple, and rational isolating intervals of width below $4\times10^{-31}$ are computed. In $[-3,1]$, $P_2$ has
one root, $x^*$, and $P_1$ has two, $r_{\mathrm{next}}$ and one near $+0.19$ (resp.\ $+0.13$). The enclosures of
$\partial_xP_2$ and of $P_1$ on $x^*\pm10^{-6}$ in Table~\ref{tab:slice} are bounded away from zero.
\emph{(E4)} Interval Cholesky~\cite{alefeldmayer1993} of $B_{4,0}(-3)-\lambda_4(-3)$ succeeds (smallest pivot lower
bounds $30.09$ and $30.40$), so $m(-3)>0$. Between $-3$ and $r_{\mathrm{next}}$ the function $\det(B_{4,0}-\lambda_4)$
vanishes only at $x^*$, and there simply. By the analytic dependence of the eigenvalue branches on $x$ (Rellich),
exactly one branch crosses $\lambda_4$ at $x^*$, and it is the lowest. A double-precision check gives
$m(x^*+10^{-3})=-0.0110$ on both slices.
\emph{(E5)} The $116$ other blocks of each slice are certified positive definite after subtracting $\lambda_j$, by
interval Cholesky on $[-3,x^*+0.02]$ with at most $14$ pieces and bisection depth $2$, at $96$ bits in
\textsc{mpmath}~\cite{mpmath}. The smallest pivot lower bound is a bound, not an eigenvalue gap; the smallest
double-precision gap is $0.051$ ($y=\sqrt3$, in the block $(1,-3/2)$) and $0.121$ ($y=3/2$)
(\path{scripts/edge_wp3_slice_iff.py}, receipts \path{wp3_slice_iff.json}, \path{wp3_slice_iff.md},
\path{wp3_slice_iff_log.txt}; about $25$\,s, and a rerun gave an identical receipt, apart from the run time, in
$24.85$\,s).
\end{proof}

The theorem is conditional: the closed-form blocks are checked against the full sector spectra to $1.5\times10^{-15}$
at $12$ random points, not re-derived symbolically inside this proof. It concerns $H+GG^*$ only. The gauge-orthogonal
edge on these slices is higher, at $-0.3552$ and $-0.3516$, and is not treated. The sectors $j>12$ are open there; in
$13\le j\le40$ every block keeps a gap of at least $5.50$ (double precision).

\subsection{Large sectors close to the edge}

\begin{theorem}[Large sectors on slabs]\label{thm:slabs}
On each box of Table~\ref{tab:slabs}, for every $(x,y,t)$ in the box and every sector $j\ge40$, every charge block of
the sector other than the twist and its mirror has lowest eigenvalue greater than $\lambda_j$. On
$[-3,-0.6]\times[0.5,3]\times[3,10]$ the same holds for every $j\ge20$.
\end{theorem}

\begin{table}[t]\centering\small
\caption{Boxes $x\times y\times t$ of Theorem~\ref{thm:slabs}, all sectors $j\ge2L_0$. $\delta$: how far the top of
the box lies below the smallest value of $\min(x^*,0)$ over the grid points of its slab (computed). Pieces: bisection
pieces for the hardest blocks. Receipts \texttt{research/edge-theorem-20261121/}\allowbreak\texttt{wp2\_interval\_proof\_*.json}.}
\label{tab:slabs}
\begin{tabular}{lcccr}\toprule
box & $L_0$ & $\delta$ & pieces & seconds\\ \midrule
$[-3,-0.575]\times[0.5,3]\times[3,4]$ & 20 & $0.1006$ & $498$ & $23.1$\\
$[-3,-0.495]\times[0.5,3]\times[3,4]$ & 20 & $0.0206$ & $551$ & $84.7$\\
$[-3,-0.355]\times[0.5,3]\times[4,5]$ & 20 & $0.1010$ & $359$ & $19.1$\\
$[-3,-0.274]\times[0.5,3]\times[4,5]$ & 20 & $0.0200$ & $384$ & $63.3$\\
$[-3,-0.229]\times[0.5,3]\times[5,6]$ & 20 & $0.1007$ & $404$ & $21.0$\\
$[-3,-0.149]\times[0.5,3]\times[5,6]$ & 20 & $0.0207$ & $954$ & $21.1$\\
$[-3,-0.183]\times[0.5,3]\times[6,7]$ & 20 & $0.1002$ & $590$ & $24.4$\\
$[-3,-0.103]\times[0.5,3]\times[6,7]$ & 20 & $0.0202$ & $1{,}579$ & $31.6$\\
$[-3,-0.154]\times[0.5,3]\times[7,10]$ & 20 & $0.1005$ & $2{,}892$ & $55.7$\\
$[-3,-0.074]\times[0.5,3]\times[7,10]$ & 20 & $0.0205$ & $12{,}445$ & $212.0$\\
$[-3,-0.02]\times[1,3]\times[7,10]$ & 20 & $0.0196$ & $126{,}082$ & $4{,}445.8$\\
$[-3,-0.6]\times[0.5,3]\times[3,10]$ & 10 & $0.1256$ & $4{,}005$ & $109.5$\\
$[-3,-0.6]\times[0.5,3]\times[3,10]$ & 20 & $0.1256$ & $1{,}653$ & $40.5$\\
\bottomrule\end{tabular}
\end{table}

\begin{proof}
The mirror block of charge $-(l+2)$ has eigenvalue exactly $\lambda_j$ and is excluded. The blocks split
into two kinds. \emph{Near-end blocks} ($l-q\le N_1$ or $l+q\le N_1$, each with at most nine states): the block is scaled to
$M=(B-\lambda_j)/l$ and written in $U=1/l$; $M$ is proved positive definite for all $U\in(0,1/L_0]$ by interval
$LDL^\top$ with bisection. \emph{Bulk rows} ($\min(\xi,2l-\xi)\ge N_1-1$): plain Gershgorin
discs~\cite{varga2001} with the quantity $G$ of Theorem~\ref{thm:domain}. All arithmetic is \textsc{mpmath}
interval arithmetic at $80$ bits~\cite{mpmath}; $N_1=8$ and $X_A=20$ for $L_0=20$, and $N_1=X_A=10$ for $L_0=10$. Every
receipt reports \texttt{all\_proved} true and the mirror exclusion checked symbolically. The near-end matrices agree
with the blocks to $4.5\times10^{-16}$ and the closed-form $G$ bounds the actual Gershgorin quantity with excess at
most $2.8\times10^{-13}$ (both double precision). A rerun of the box $x\le-0.6$, $L_0=20$ reproduced the proof in
$125.7$\,s (\path{scripts/edge_wp2_interval_proof.py}, receipts as in the caption and
\path{wp2_verify_rerun_x-0.6_L20.json}; design in \path{wp2_large_sector_bound.md}).
\end{proof}

Four limits. The column $\delta$ is measured from the worst grid point of each slab, so between grid points the
distance to the edge is not proved. The sectors $j<40$ on these boxes are not done, so the theorem does not by itself
make the twist the floor there. And $x\in(-0.02,0)$ is not covered. Finally, the scaled near-end matrices and the
bound $G$ are checked against the blocks numerically (the two double-precision checks above), not as symbolic
identities; the theorem rests on those checks. (The decimal box endpoints are entered as doubles; for instance the
top $-0.02$ is the double just below $-0.02$, a negligible difference.)

\begin{remark}[Why plain Gershgorin is not enough near the edge]\label{rem:gersh}
At the near end of a sector ($\xi=0$) the diagonal gap is $lK_a$, while the ladder coupling grows like $\sqrt l$. So
plain Gershgorin needs $l$ of order $1/K_a^2$, and $K_a$ becomes small as $x\to0$. The true margin is $lK_a+O(1)$.
Sampled cutoffs (double precision; not certified) show the cost: at distance $0.1$ below the edge $\min(x^*,0)$, the
worst grid point $(t,y)=(10,3)$ needs $l\approx8978$ for plain Gershgorin and $l\approx353.5$ for a weighted
Gershgorin bound with an $l$-dependent diagonal similarity; a direct check puts the true weighted cutoff there at $345$.
The cutoffs are grid-quantised estimates. The weighted bound was not turned into an interval proof; the near-end
$LDL^\top$ of Theorem~\ref{thm:slabs} replaced it (\path{wp2_bounds.json}).
\end{remark}

\subsection{The all-sector edge}

\begin{conjecture}\label{conj:edge}
For $t\ge3$ and $y>0$, the all-sector edge of $H+GG^*$ equals $\min(x^*_{12}(t,y),0)$.
\end{conjecture}

The evidence for it lies in $3\le t\le12$, $0.25\le y\le3$: the competitor map, the tail checks above, and the
near-top families below.

\paragraph{The conjecture fails near $t=2$ (double precision; two builders).} The same statement on $2<t\le12$ is
false. At $(t,y)=(2.01,0.25)$, $x^*_{12}=-0.62410$. At $x=-0.62460$, where $\phi_A$ over $j\le12$ is still $+0.0052$,
the block $q=8$ of the sector $j=16$ lies below $\lambda_{16}$ by $0.00492$, both from the blocks and from the general
$H+GG^*$ builder (two eigenvalues below), and by $0.00353$ in the gauge-orthogonal spectrum. At $(2.001,0.25)$, at
$x=-0.64263$ and $x=-0.64513$, both below $x^*_{12}=-0.64213$, the sector $j=58$ lies below by $0.2387$ and
$0.2109$. A scan of $j\le3000$ at $t=2.001$, $2.01$ and $2.03$, $y=0.25$ and $0.5$, found an undercut only at
$y=0.25$ for $t=2.001$ and $2.01$ (\path{research/edge-theorem-20261121/workflow_reports.json}, the competitor-map
rerun).

\paragraph{The mechanism: two near-top families.} The blocks $q=l+1$ (three states) and $q=l$ (six states) can be
written with $l$ symbolic. Their determinant numerators are $D_+(x,y,s,l)$, of degrees $(4,4,9,3)$, and $D_0(x,y,s,l)$, of
degrees $(8,10,19,6)$, with $s=\sqrt t$ (checked numerically against the blocks to $6.8\times10^{-13}$). Their leading
coefficients in $l$ are
\[
  c_+=-128\,s\,x\,(s-1)(s+1)(s^2-2),
\]
\[
  c_0=-65536\,s\,x\,(s-1)(s+1)(s^2-2)(3s^2-4)(sx-2s^2+2)(sx-s^2+2),
\]
and both vanish identically at $s=\sqrt2$, that is at $t=2$. Away from $t=2$ the factor $x$ pulls the family edges
towards $x=0$ as $j$ grows. The factor $(sx-s^2+2)$ of $c_0$ vanishes at $x=(t-2)/\sqrt t$ ($0.5774$ at $t=3$), so the
statement that $x=0$ is the only root of $c_0$ in the range of the edge holds only for $x\lesssim0.5$. These polynomials
and coefficients are exact algebra (\path{scripts/edge_wp1_block_polynomials.py}, receipts \path{wp1_top_family.txt},
\path{wp1_family_members.json}); the family edges below are double precision. At $(t,y)=(12,2)$ the edge of the family
$q=l+1$ falls roughly like $1/j$ (Table~\ref{tab:family}). At $t=2$ the family edges appear instead to converge to a
negative limit: the family $q=l$ at $j=2000$ gives $-0.68467$ for $y=0.25$ and $-0.61843$ for $y=0.5$.

\begin{table}[t]\centering\small
\caption{Edge of the near-top family $q=l+1$ at $(t,y)=(12,2)$ (double precision; \texttt{wp1\_competitor\_map.md},
Section~4).}\label{tab:family}
\begin{tabular}{lcccccccc}\toprule
$j$ & $12$ & $24$ & $48$ & $96$ & $200$ & $400$ & $1000$ & $2000$\\ \midrule
edge & $0.2075$ & $0.1458$ & $0.0905$ & $0.0513$ & $0.0264$ & $0.0137$ & $0.0056$ & $0.0028$\\
\bottomrule\end{tabular}
\end{table}

Six of the eight competitors found in $j\le12$ belong to these two families: $D_+$ at $l=\tfrac12,1,6$ gives
$(1,3/2)$, $(2,2)$ and $(12,7)$; $D_0$ at $l=1,6$ gives $(2,1)$ and $(12,6)$; $D_0$ at $l=\tfrac12$ gives $(1,1/2)$ up to a
spurious factor. Only $(4,0)$ and $(3,1/2)$ are sporadic. Because $c_+$ and $c_0$ vanish at $t=2$, near $t=2$ the
families need not fall back to $x=0$, and there they can take over from the $j\le12$ competitors. That is what the
refutation shows, and it is why the conjecture is stated for $t\ge3$ only.

\section{The floor on a domain}\label{sec:domain}

Inside its region the twist can be proved to be the floor of every sector on a box of positive volume.

\begin{theorem}[The floor on a domain]\label{thm:domain}
On $D=\{x\in[-3,-1.5],\ y\in[0.5,3],\ t\in[3,10]\}$ the twist is the lowest eigenvalue of $H+GG^*$ in every sector,
with multiplicity exactly two, and the lowest gauge-orthogonal eigenvalue there.
\end{theorem}

It is also Theorem~7.3 of the version~3 draft of~\cite{davisfloor2026}; it is not in the published version~2.

\begin{proof}
\emph{(a) Reduction to charge blocks.} By Lemma~\ref{lem:blocks} each sector is the orthogonal sum of its charge
blocks. The two twist blocks are $1\times1$ with entry $\lambda_j$ (Proposition~\ref{prop:closed} and the
antiunitary map $J\otimes K$ that reverses the charge). It suffices to show that every other block has lowest
eigenvalue above $\lambda_j$ on $D$. Then $\lambda_j$ is the lowest eigenvalue of $H+GG^*$ with multiplicity exactly two,
and since the twist is gauge-orthogonal and $H+GG^*=H$ on the slice, it is also the lowest gauge-orthogonal
eigenvalue. The closed-form rows, built symbolically with $x$ and $y$ free, agree with the general sector builder to
$1.7\times10^{-13}$ (receipt entry \texttt{point\_check}).

\emph{(b) Sectors $j\ge80$.} For the row $m=l-\xi$ of type $a$ the Gershgorin quantity is
\[
  G=4\xi(2l-\xi)\Bigl(1-\frac1t\Bigr)+lK_a+E_a-(\text{ladder}),\qquad K_a=4-\frac8t-\frac{4(w_a+x)}{\sqrt t},
\]
built on the identity in the proof of Theorem~\ref{thm:xpos}. On $D$, $K_a>0$ for every non-twist type. $G$ is
proved positive for all $l\ge40$ on $D$ by interval bisection, split into region~A ($\xi\le20$) and region~B
($\xi\ge20$, with a relaxed bound). In region~B the smallest lower bounds run from $0.7226$ (type $(0,0)$) to $1.5918$
(type $(-1,1)$). In region~A the smallest lower bound is $1.0678\times10^{-4}$ (type $(0,1)$, $\xi=0$, $1360$ pieces); this
is the overall minimum. The closed-form $G$ is a lower bound on the actual Gershgorin quantity, checked on random rows
with excess at most $2.84\times10^{-13}$ (double precision). The run takes $13.5$\,s
(\path{scripts/ricci_su2_family_large_j.py}, receipt \path{family_large_j.json}).

\emph{(c) Sectors $j<80$.} Every non-twist block is certified above $\lambda_j$ on $D$ by a mean-value interval
certificate. $D$ is cut into boxes. On each box the lowest eigenvalue of $B-\lambda_jI$ at the centre is bounded below
by an interval $LDL^\top$ certificate with outward rounding and an a~posteriori residual
(\path{scripts/ricci_su2_outward.py}). Its variation over the box is bounded by
$\sum_i\sup\lVert\partial_iA\rVert_F$ times the half-widths, with the derivatives enclosed in interval
arithmetic~\cite{mpmath}. A box whose bound is not positive is split. In all, $80$ sectors and $3400$ blocks were
certified on $1{,}640{,}935$ boxes with no failure. The smallest certified margin is $1.2220726\times10^{-6}$, in the
sector $j=1$ (Table~\ref{tab:domain}). The summed run time is $229{,}134.7$\,s, about $63.6$ CPU-hours
(\path{scripts/ricci_su2_family_block_cert.py}, receipt \path{family_block_cert_j80.json}, \texttt{all\_certified}
true). Receipts are in \path{research/su2-floor-family-20261119/}.
\end{proof}

\begin{table}[t]\centering\small
\caption{The small-sector certificate of Theorem~\ref{thm:domain}: the sector $j=0$, the four sectors with the
smallest margins, and the total (\texttt{family\_block\_cert\_j80.json}).}\label{tab:domain}
\begin{tabular}{lrrr}\toprule
sector & boxes & smallest margin & seconds\\ \midrule
$j=0$ & $11{,}431$ & $1.41\times10^{-4}$ & $400.0$\\
$j=1$ & $71{,}052$ & $1.2220726\times10^{-6}$ & $3{,}152.2$\\
$j=2$ & $71{,}729$ & $2.94\times10^{-5}$ & $5{,}360.8$\\
$j=3$ & $54{,}736$ & $5.77\times10^{-6}$ & $5{,}501.7$\\
$j=36$ & $14{,}015$ & $6.90\times10^{-5}$ & $2{,}103.3$\\ \midrule
$j=0,\dots,79$ & $1{,}640{,}935$ & $1.2220726\times10^{-6}$ & $229{,}134.7$\\
\bottomrule\end{tabular}
\end{table}

\begin{remark}[Why $D$ stops at $x=-1.5$]\label{rem:Dstop}
Plain Gershgorin needs large $l$ unless the domain stays away from the edge (Remark~\ref{rem:gersh}). Sampled cutoffs
(double precision) are $l\approx640$ for $x\in[-2,-0.5]$, $l\approx160$ for $[-2,-1]$ and $l\approx40$ for $[-3,-1.5]$. The
near-end treatment of Theorem~\ref{thm:slabs} is the route closer to the edge; it covers $j\ge40$ up to about $0.02$
below the grid edge, but the small sectors there remain to be certified.
\end{remark}

\begin{remark}[$D$ lies inside the region]\label{rem:Dinside}
On $D$, $x\le-1.5<0$, and $r_{\mathrm{lo}}<-7.33$ on all of $D$, so $D$ sits well inside the bounds of
Corollary~\ref{cor:interval}. In double precision, $\min\phi_A$ at $x\in\{-3,-1.5\}$ over an $11\times15$ grid of $D$'s
$(y,t)$ range is $2.35$ ($j\le12$) and $\min\phi_S$ on a $3\times3$ grid is $2.69$. These grid values are evidence only;
the inclusion itself is settled by the certificate of Theorem~\ref{thm:domain}.
\end{remark}

\section{What decides the edge, and what does not}\label{sec:decide}

Three simpler readings of why the twist is the floor were tested (Tests~2, 4 and~5), each written down before its
data, and each failed. A fourth test (Test~6), of a two-parameter model for the edge, failed in both of its
pre-registered predictions.

\paragraph{Symmetry (Test~2; failed).} The reading was that the rotation symmetry itself makes the maximal-charge state
the floor. On the squashed meron, which has the same symmetry, the twist is an exact gauge-orthogonal eigenvector, yet it
is the floor of no sector at any of the $2{,}600$ pairs $(t,j)$ checked (eigenvector residual at most
$4.7\times10^{-12}$, $\lVert G^*b\rVert\le8.5\times10^{-15}$). The Berger field, run through the same code, became the
floor of every sector $j\le12$ from $t=2.70$. In two $SU(3)$ families with the Maurer--Cartan field $A=\tfrac12\theta$,
one stretching the $U(2)$ fibre and one the hypercharge circle, the extremal state is the floor of none of nine sectors
at any of $200$ grid points; one of the $1{,}800$ point-sector pairs needed a fallback singular value routine. One
difference, observed after the test and not pre-registered, is in the curvature term
$R_0=\mathrm{Ric}\otimes I-2[F,\cdot\,]$: on the twist's internal frame the field-strength share is $8-6t$ on the Berger
field, but positive on the squashed meron, from $+1.26$ at $t=2$ to $+0.40$ at $t=8$. Tests~4 and~5 then showed that
this term does not decide the floor (\path{scripts/ricci_su2_test2_family.py}, \path{scripts/ricci_su2_gyromagnetic_split.py},
receipts \path{test2_meron.json}, \path{su3_test2_family2a.json}, \path{su3_test2_family2b.json}, \path{gyromagnetic_split.txt}).

\paragraph{Alignment (Test~4; failed).} The reading was that a field strength aligned with the twist's spin makes it the
floor. For $A=s\theta$ on the round sphere, $s=1.05,\dots,3.00$ ($40$ points), the field strength is aligned throughout
(its field-strength share $s_F$ on the twist's frame runs from $-0.42$ at $s=1.05$ to $-48$ at $s=3$). In $SU(2)$ the
twist is the floor of $j=0$ for $s\ge1.5$ and of no sector $j\ge1$ at any $s$; in the sector $j=1$ the gap grows from
$16.5$ to $35.9$. In $SU(3)$ it is the floor of the sector $(0,0)$ for $s\ge1.5$ and of no other
(\path{scripts/ricci_alignment_test.py}, receipts \path{alignment_test_4a.json}, \path{alignment_test_4b.json}).

\paragraph{Alignment and the curvature term (Test~5; failed).} The reading was that the twist is the floor when it alone
spans the lowest eigenspace of $R_0$. On $A=s\theta$ over $(t,1,1)$, $72$ points were classified from $R_0$ alone
before any floor was computed: $35$ singled out, $5$ shared, $32$ not lowest. The twist is the floor of every sector at
none of them. At every singled-out point it is the floor of $j=0$ only. The gap grows with $j$ (at $s=1.5$, $t=3$:
$6.6$, $12.6$, $25.7$, $41.1$ for $j=1,2,4,8$) and shrinks with the stretch (at $j=1$, $s=1.5$: $13.8$, $6.6$, $2.1$ at
$t=1.5$, $3$, $8$) (\path{scripts/ricci_stretch_test.py}, receipt \path{stretch_test.json}).

\paragraph{A model for the edge (Test~6; failed).} The model $x^*(t)=c_1/t+c_2/t^2$ was fitted to the gauge-orthogonal
edge of $j\le12$ at $t=2,3,4,6$ and heights $y=0.5$, $1.5$, $\sqrt t$. Prediction~E1 asked for agreement to $0.05$ at
held-out stretches. Interpolation held at $6$ of $6$ points ($t=2.5$ and $5$; errors $-0.018$, $-0.038$, $-0.047$,
$+0.010$, $+0.024$, $+0.029$). Extrapolation to $t=8$ held at $y=0.5$ ($+0.033$) but failed at $y=1.5$ (predicted
$-0.052$, measured $+0.040$, error $+0.092$) and at $y=\sqrt t$ (predicted $-0.041$, measured $+0.138$, error $+0.179$).
Prediction~E2 failed: the model puts the Berger crossing at $t_c=2.6315$, against $t^*=2.6980$, a distance
$0.067>0.05$ (\path{scripts/ricci_su2_edge_test.py}, receipts \path{edge_test_predictions.json},
\path{edge_test_results.json}). Two later findings bear on this test. The positive edges measured at $t=8$ are
artefacts of the cut-off $j\le12$ (Section~\ref{sec:interval}), and the region reported there as $[-3,x^*]$ is part of
an interval bounded below (Proposition~\ref{prop:lower}).

\paragraph{Neutrality.} A further reading is that the competitor is neutral, with total fibre charge zero, as the
lowest state of $B_{4,0}$ is. Data already in hand count against it as a necessary condition. At $y=0.5$ the
competitors are blocks of non-zero charge (Section~\ref{sec:edge}), and at $14$ grid points with $y\ge1.1$ and
$x^*_{12}\le0$ the competitor is $(1,\pm3/2)$ or $(2,\pm2)$ (\path{wp1_competitor_map.json}). This finding predates
pre-registration~B (Section~\ref{sec:prereg}) and is not one of its results.

\paragraph{What does decide it.} In each sector, whether the twist is the floor is a comparison between $\lambda_j$ and
the lowest eigenvalues of that sector's other charge blocks. Its structure is exact. In every block, $x$ enters the
diagonal through $(w_a-2m/\sqrt t)^2$, with $w_a$ affine in $x$; every entry of the remaining matrix $C$ has degree at
most one in $x$; and the ladder couplings do not involve $x$ (Section~\ref{sec:setting}). In the twist, $x$ enters
twice: by minimal coupling, through the fibre momentum $\kappa$, and by spin, through the field strength
$B=y^2-2\sqrt t\,x$ along the fibre; the slope is $\partial_x\lambda=-2\kappa+4\sqrt t$ (Section~\ref{sec:twist}).
Numerically the two sides behave very differently near the edge. At three points, the Berger field at $t^*$,
$(t,y,x)=(4,2,-0.19)$ and $(8,\sqrt8,0.14)$ (the last at a positive $j\le12$ edge, itself a cut-off artefact), the
twist's slope in the sector $j=4$ is $9.741$, $9.620$ and $10.179$, while the lowest eigenvalue of $B_{4,0}$ has slope
$-1.561$, $-0.623$ and $+1.684$ (double precision; \path{twist_anatomy.json}). The competitor is nearly blind to $x$;
the twist is not. An exploratory map, not pre-registered, agrees: over $t=1,2,3,4,6$, $x\in[-3,3]$ ($49$ values),
$y\in[0,3]$ ($37$ values) and $j\le8$, all $3{,}294$ cells where the twist is the floor of all nine sectors have
$x\le0$, and $3{,}286$ have $x<0$ (\path{scripts/ricci_su2_twist_family_map.py}, receipt \path{twist_family_map.json}).
This is consistent with Theorem~\ref{thm:xpos}.

\paragraph{The literature.} The twist is a cousin of the unstable mode of a charged vector field in a homogeneous
magnetic background~\cite{nielsenolesen1978}, the classical face of the chromomagnetic instability whose one-loop form
appears in~\cite{savvidy1977}; see also~\cite{tanjiitakura2012} for the standard usage. The flat-space mechanism is
stated explicitly in~\cite{bcgk2014} as a Landau-level energy plus a Zeeman term with gyromagnetic factor~$2$. The
decomposition of $\lambda$ in Section~\ref{sec:twist} has the same shape; I note the resemblance only, and no limit
relating the two is proved. Instabilities of homogeneous backgrounds on the round $S^3$ are treated
in~\cite{avramidicollopy2012,klc2021}, and spectra on Berger spheres in~\cite{lauret2019,henkellauret2026}. The closest
is~\cite{avramidicollopy2012}, on the negative modes of covariantly constant backgrounds, on the round sphere only. In
the literature checked I found no edge of this kind and no Berger-sphere family of this kind; several classical full
texts could not be read, so this is a statement about the search, not a priority claim.

\section{Pre-registered tests A, B and C}\label{sec:prereg}

Three further tests were written down, reviewed and fixed before any scored data were computed, as were Tests~2, 4, 5
and~6 of Section~\ref{sec:decide}. This section reports their designs, gates, amendments and results. Each amendment was
made after a code gate fired (for Amendment~C1, after a smoke test showed that a gate would fire) and before any scored
quantity was computed. Their documents are
\path{research/edge-theorem-20261121/wp4_prereg_A_heldout.md} (with \path{wp4_prereg_A_points.json}),
\path{wp4_prereg_B_general_fields.md} and \path{wp4_prereg_C_su3.md}.

\subsection{Test A: the edge at held-out points of the family}

\paragraph{Points.} $48$ points $(t,y)$, drawn with seed $20261121$: $30$ in the class where $P_2$ governs
($t\sim U[2.2,11.5]$, $y\sim U[1.1,3.2]$), $10$ at small $y$ ($y\sim U[0.3,0.7]$), and $8$ on a fixed grid
($t\in\{2.2,3.5,7.0,11.0\}$, $y\in\{1.2,2.5\}$). None lies within $10^{-3}$ of the $50$ pairs $(t,y)$ used before.

\paragraph{Predictions.} A1: the edge is a root of $P_2$, with competitor $B_{4,0}$, at the $38$ points outside the
small-$y$ class. A2: at small $y$ the competitor is a block of the sector $j=2$ or $3$. A3: no sector $13\le j\le79$
undercuts the twist at three values of $x<0$ (just below $\min(x^*,0)$, at the midpoint of the region, and just above
$x_{\mathrm{lo}}$); $x>0$ is excluded, since the undercut there is proved. A4: on a $0.01$ grid over
$[r_{\mathrm{lo}}-0.5,\,x^*+0.5]$ the region of the sectors $j\le12$ is a single interval.

\paragraph{Measurement and gates.} The edge is measured on an independent code path, the full $H+GG^*$ sector, with
both copies of $\lambda_j$ removed. Gates: A-code-1 (blocks against the full spectrum, $10^{-12}$ relative); A-code-2
($18$ known rows to $10^{-9}$); A-code-3 (the Berger path at $t_A$, $10^{-8}$ relative); A-mult (exactly two eigenvalues
at $\lambda_j$); A-branch (the largest real root of $P_1$ must miss the measured edge by more than $10^{-6}$ at no fewer
than $80\%$ of the $P_2$-class points, or A1 is inconclusive); A-bound ($x_{\mathrm{lo}}\ge r_{\mathrm{lo}}$); A-bound-hi
($x^*\le r_{\mathrm{hi}}$). Verdict rule: A1 is supported if it holds at all $38$ points, and a failure counts against the
edge clause; a failure of A3 counts against the claim for $x<0$; a hole in A4 counts against the interval. Three caveats
were written down in advance. A1 agreement is partly automatic, because $P_2$ comes from the same closed-form rows that
A-code-1 checks. ``Held-out'' means not yet computed. The claims cover $j\le79$ only.

\paragraph{Amendment~1.} Gate A-code-3 passed (relative residual $1.0\times10^{-18}$). Gate A-code-2 fired, with
maximum deviation $3.67$. On the $18$ known rows alone, $P_2$ was recovered exactly (degrees $(7,8,16)$, $113$ terms) and the
edge is a real root at every row, but not the largest root in $[-3,3]$. $P_2$ has three or four further real roots above
the edge, from about $0.55$ at $t=2$ to $3.5$ at $t=12$, $y=\sqrt t$, where a higher eigenvalue of $B_{4,0}$ meets
$\lambda_4$. The rule was amended: the predicted edge is the largest real root $r$ of $P_2$ in $[-3,3]$ at which the lowest
eigenvalue of $B_{4,0}(r,y,t)$ equals $\lambda_4(r)$ to $10^{-8}\max(1,|\lambda_4|)$; if there is none, A1 fails there.
A-branch uses the same selection for $P_1$. Nothing else changed. The known-answer rows thus count against the
statement ``the edge is the largest real root of $P_2$''.

\subsection{Test B: general circle-symmetric fields}

\paragraph{B-I, a code control.} The class of Proposition~\ref{prop:general} reduces to the ansatz by constant gauge
transformations. B-I checks this in code and is not evidence. $200$ members are drawn (generator seed $1122$, weight
$k=\pm1$, $x\sim U[-10,10]$, $y\sim U[0,6]$, $z\sim U[-6,6]$, $t\sim U[0.5,15]$, a random constant gauge frame), sectors
$j\le8$. Gates: B-sym, B-exact, B-red (the twist's value equals $\lambda$ at the reduced coordinates to $10^{-10}$
relative), B-spec (full spectra equal those of the reduced field to $10^{-12}$ relative) and B-even ($10$ extra draws at
$y=z=0$).

\paragraph{B-II, the part that can fail.} A random scan of the reduced family over $(\rho,t)$, $\rho=\sqrt{y^2+z^2}$,
at the $200$ reduced draws, with $x$ and $k$ discarded and the edge located on $x\in[-10,10]$ in $j\le12$. Gates:
B-II-code (the $27$ rows of \path{exact_edge.json} reproduced to $10^{-9}$); B-count (at least $60$ scored points, or B-S
and B-half are inconclusive); B-S (the competitor's slope ratio $|\partial_x\lambda_{\mathrm{comp}}|/|\partial_x\lambda|<0.5$
at every scored edge; the threshold was chosen after seeing ratios $0.16$, $0.064$ and $0.165$, so a pass is weak
evidence); B-half (exactly two sign changes of $\phi_A$ on a $0.01$ grid over $[r_{\mathrm{lo}}-0.5,\,x^*+0.5]$);
B-N, the competitor's charge, descriptive only.

\paragraph{Amendment~B1.} Gate B-II-code stopped the run before any B-II point was computed. At all $27$ known rows the
edge was reproduced (largest deviation about $4\times10^{-15}$), but at two rows ($t=3$ and $t=4$, height $0.5$) the
competitor was read as $(2,+1)$ where the receipt has $(2,-1)$, and the reverse. At every known row whose competitor has
$q\ne0$, the blocks $(j,q)$ and $(j,-q)$ have the same lowest eigenvalue to $6.7\times10^{-15}$ (double precision), and which
one is reported is decided by rounding. The gate was amended to accept the competitor up to $(j,q)\leftrightarrow(j,-q)$
when the two lowest eigenvalues agree to $10^{-12}$ relative at $x^*+10^{-4}$. The scored arms are unaffected: B-N asks
whether $q=0$, B-S uses a slope that the mirror shares, and B-half does not use the competitor.

\subsection{Test C: $SU(3)$ with a free fibre field}

\paragraph{The family.} On $SU(3)$ with the metric that stretches the hypercharge circle by $t$ (on $\lambda_8$), the
left-invariant fields commuting with the joint frame-and-gauge action of $U(2)$ form a space of dimension $4$ (computed
by SVD to machine precision). Up to gauge they are $M_a=uX_a$ ($a=1,2,3$), $M_8=xX_8$, $M_a=yX_a$ ($a=4,\dots,7$),
$y\ge0$, with components in the orthonormal coframe and $x$ the field along the stretched fibre. The Maurer--Cartan
member $u=y=\tfrac12$, $x=1/(2\sqrt t)$ is the field of the second $SU(3)$ family of Section~\ref{sec:decide}, where the
extremal state was the floor nowhere; here $x$ is free. The analogue of the twist has charge $q_8=\sqrt3/2$ under the
fibre generator. Three structural statements were labelled conjectures until their gates pass: F1, the coefficient of
$x^2$ in its eigenvalue is $q_8^2=3/4$ in every sector; F2, its region in $x$ is bounded, via the gauge-orthogonal state
$\hat\theta^8\otimes X_8$ of the sector $(0,0)$, whose Rayleigh quotient is affine in $x$; F3, the lift discount is linear
in $x$ with one sign $\sigma$ across the sectors.

\paragraph{Predictions.} C-P1, existence: with $x$ free, the region where the twist is the floor of all nine sectors
(gauge-orthogonal spectrum) is non-empty, strictly at every grid point with $t\ge3$, weakly at one point at least.
C-P2, side: at every informative point the midpoint of the region satisfies $\sigma x<0$. C-P3: the region is one
interval. C-P4: at the upper edge the competitor's slope ratio is below $0.5$ (scored); its charge is descriptive.
C-P5: the upper edge against $t$, descriptive. Every claim of Test~C is about the nine listed sectors.

\paragraph{Grid and gates.} $t\in\{2,3,4,6,8\}$ and $(u,y)\in\{(\tfrac12,\tfrac12),(\tfrac12,1),(1,1),(1,2)\}$, $20$ shape points,
the nine sectors $(p,q)$ of Test~2 listed below, $x$ scanned in steps of $0.25$ on $[r_{\mathrm{lo}}-0.5,\,r_{\mathrm{hi}}+0.5]$. Gates:
C-K1 to C-K5 (the earlier $SU(3)$ controls, an $SU(2)$ analogue against \path{exact_edge.json}, and a known answer at
the Maurer--Cartan member), C-sym, C-quad (second difference $3/2$ to $10^{-9}$), C-lin, C-sigma, C-mult and C-bound.

\paragraph{Amendment~C1.} A smoke test of the structure gates at one shape point, $(t,u,y)=(8,1,2)$, was run before any
floor, edge or competitor was computed. C-quad passed ($3/2$ to $5.3\times10^{-15}$ in all nine sectors), C-lin passed
(second difference $0.0$), $\sigma=+1$, and $[r_{\mathrm{lo}},r_{\mathrm{hi}}]=[-7.656,4.120]$. Gate C-mult would have stopped
the run: the number $n_b$ of slice eigenvalues equal to the twist's value was $6$, $4$, $8$, $5$, $5$, $6$, $10$, $6$, $7$ in the sectors
$(0,0)$, $(1,0)$, $(1,1)$, $(2,0)$, $(2,1)$, $(3,0)$, $(2,2)$, $(3,1)$, $(4,0)$, not $1$ or $2$. The family commutes with a non-abelian $U(2)$, so
eigenvalues come in multiplets of its $SU(2)$; the twist has spin $p/2+1$ there, and the count is
$n_b=(p+3)(2\text{ if }p=q\text{, else }1)$, which matches all nine. This formula was derived after seeing the counts.
The amended gate requires that count and removes exactly $n_b$ copies. The predictions C-P1 to C-P5 are unchanged. These
gate values are design-stage quantities and say nothing about the predictions.

\paragraph{Amendment~C2.} The rerun passed C-K1 to C-K5 again (C-K5 to $2.6\times10^{-14}$), and the shape point
$(2,\tfrac12,\tfrac12)$ passed every structure gate. At $(2,\tfrac12,1)$ the amended C-mult stopped the run: the counts
were $6$ in $(1,0)$, where the formula gives $4$, and $12$ in $(1,1)$, where it gives $8$, the same at all three $x$
values. The extra states form a second $SU(2)$ doublet, with charges $3$ and $4$, outside the twist's multiplet, whose
eigenvalue equals the twist's identically in $x$ at $(u,y)=(\tfrac12,1)$; this was seen at $t=2$, $2.0001$ and $3$, and
not at the other shapes examined. They are tied competitors, not copies of the twist. The gate was amended again: $n_b$
must be at least the multiplet count $(p+3)(2\text{ if }p=q\text{, else }1)$ and the same at the three $x$ values, exactly
the multiplet count is removed, and any further states at the twist's value stay in the spectrum as competitors. Under
the registered non-strict floor rule a tie leaves the twist the floor; the extra counts are reported, and where they are
identical in $x$ the twist is never the unique floor in that sector. In the same amendment the spectra were moved, as an
implementation change, to the eigenbasis of the field-independent charge, block by block (a unitary change of basis,
agreeing with the dense computation to $2\times10^{-13}$ on the largest sector). The predictions are unchanged, and no
floor, edge or competitor had been computed.

\subsection{Results}

\paragraph{Test A (COMPUTED).} Receipt: \path{research/edge-theorem-20261121/wp4_heldout.json}. Every gate passed:
A-code-1 to $4.1\times10^{-15}$; A-code-2 with $113$ terms and the $18$ known rows to $2.4\times10^{-13}$ under
Amendment~1; A-code-3 to $1.0\times10^{-18}$; A-mult at every point and sector $j\le12$; A-bound and A-bound-hi, with
$x_{\mathrm{lo}}-r_{\mathrm{lo}}\ge1.61$ everywhere; A-branch at $100\%$ of the $P_2$-class points.
\begin{itemize}
  \item \emph{A1 failed at $4$ of $38$ points.} At the other $34$ the competitor is $B_{4,0}$ and the measured edge
  equals the predicted root of $P_2$ to $1.35\times10^{-12}$. The failures are listed in Table~\ref{tab:a1fail}. At three
  of them, at small stretch with a strong cross field, the block $(1,-\tfrac32)$ reaches the twist first; the map of
  Section~\ref{subsec:competitors} already showed $(1,\pm\tfrac32)$ there for $y\ge2$ and $t\le4$, and the
  $P_2$-class box of Test~A overlapped that corner. At the fourth, $(t,y)=(11,1.2)$, the block $(12,-7)$ of the cut-off
  sector sets a positive edge in $j\le12$, which Theorem~\ref{thm:xpos} already makes irrelevant to the all-sector
  question. The statement ``for $y\gtrsim1$ the edge is the root of $P_2$'' is therefore counted against; what holds is
  that $P_2$ governs where $B_{4,0}$ is the competitor.
  \item \emph{A2 held} at all $10$ small-$y$ points: $(2,\pm1)$ eight times and $(3,-\tfrac12)$ twice.
  \item \emph{A3 held} at all $48$ points: no non-twist block of any sector $13\le j\le79$ lies below the twist at the
  three values of $x<0$; the smallest margin is $+0.52$.
  \item \emph{A4 held} at all $48$ points: the region in $j\le12$ is a single interval on the $0.01$ grid.
\end{itemize}
Descriptive. The lower edge $x_{\mathrm{lo}}$ ranges over $[-7.81,-3.64]$ and is always set by the sector-$0$ block of
charge $-1$ (with its mirror). The gauge-orthogonal edge lies above the $H+GG^*$ edge at every point, by $1.5\times10^{-4}$
to $0.091$. The Test~6 model $c_1/t+c_2/t^2$ misses the measured edge by $0.076$ on average and $0.20$ at worst. Of the
$48$ upper edges in $j\le12$, $20$ are positive; these are the cut-off artefacts of Theorem~\ref{thm:xpos}.

\begin{table}[h]\centering\small
\begin{tabular}{rrrrl}\toprule
$t$ & $y$ & predicted (root of $P_2$) & measured & competitor\\\midrule
$2.295$ & $2.687$ & $-0.6089$ & $-0.6941$ & $(1,-\tfrac32)$\\
$2.2$ & $2.5$ & $-0.6378$ & $-0.7148$ & $(1,-\tfrac32)$\\
$3.5$ & $2.5$ & $-0.2979$ & $-0.3202$ & $(1,-\tfrac32)$\\
$11.0$ & $1.2$ & $+0.0847$ & $+0.0736$ & $(12,-7)$\\\bottomrule
\end{tabular}
\caption{The four points where A1 failed.}\label{tab:a1fail}
\end{table}

\paragraph{Test C (COMPUTED).} Receipt: \path{research/edge-theorem-20261121/wp4_su3_free_fibre.json}. Every gate passed:
C-K1 to C-K3; C-K4 (the $SU(2)$ analogue: coefficient of $x^2$ equal to $1$ to $3\times10^{-14}$, three known edges to
$3\times10^{-10}$); C-K5 to $2.6\times10^{-14}$; C-sym ($\norm{[Q,H]}$ to $1.1\times10^{-14}$ relative, residual of the
twist $8\times10^{-16}$, $\norm{G^*b}$ $7\times10^{-16}$); C-quad at all $20$ shapes and nine sectors to
$9\times10^{-15}$, so F1 holds; C-lin exactly, so F2 holds; C-sigma with $\sigma=+1$ everywhere; C-mult under
Amendment~C2, with the extra tied doublet at the five shapes $(u,y)=(\tfrac12,1)$.
\begin{itemize}
  \item \emph{C-P1 failed, weakly and strictly.} At none of the $20$ shapes is the twist the floor of all nine sectors
  for any $x$ on the scan, for either spectrum. By the rule fixed in advance this counts against the reading of
  Section~\ref{sec:scope} beyond $SU(2)$, in this family.
  \item C-P2 is not informative, and C-P3, C-P4 and C-P5 are vacuous, because there is no region.
\end{itemize}
A diagnostic run after the verdict, and not part of the protocol (\path{wp4_su3_diagnostic_posthoc.json}), shows how far
the twist is from the floor. At $t=4$, at the $x$ that maximises its worst margin over the sectors, the margin is negative
in every sector other than $(0,0)$. It grows more negative with the sector, from $-1.5$ to $-2.6$ in $(1,0)$ to $-6.1$
to $-8.3$ in $(4,0)$, and $(4,0)$ is the blocking sector at $121$ of the $127$ values of $x$ scanned. In $SU(2)$ the field along the
circle gives the twist a discount on the lift to higher sectors (Corollary~\ref{cor:lift}). In this $SU(3)$ family the
higher sectors cost it more at every $x$ allowed. What carries over is structural: the coefficient of $x^2$ is the
squared charge, and the bounding state is affine.

\paragraph{Test B.} B-I passed every gate at all $200$ draws: B-sym $1.7\times10^{-15}$, B-exact $3.5\times10^{-15}$,
B-red $2.2\times10^{-14}$, B-spec $3.9\times10^{-15}$, B-even $3.7\times10^{-16}$. B-II-code passed under Amendment~B1.
The B-II scan is running, and its results will be reported here.

\section{Scope and open problems}\label{sec:scope}

\paragraph{Scope.} The results concern the classical linearised operator. Both the $L^2$ (Coulomb) slice and $H+GG^*$
are given, and where only one is proved the text says which. Nothing here is a statement about the quantum theory or a
mass gap. Most members of the family are not Yang--Mills critical points, so the edge describes the landscape at a
field, not a dynamical transition; at the Berger field the Yang--Mills flow moves away faster than any twist grows
(\cite{davisfloor2026}, ``Further observations''). The stability theorems for Yang--Mills
fields~\cite{bl1981,ohnitapan1990,kot1986,stern2010} and the recent floor-type estimate for abelian Yang--Mills--Higgs
fields on the round sphere~\cite{jin2026} concern critical points. The class of fields is the one of
Proposition~\ref{prop:general}: left-invariant, symmetric under the right circle about the long axis lifted by a constant
gauge circle. Outside $D$ and the two slices of Theorem~\ref{thm:slice}, every statement about the region is computed in
$j\le12$ or bounded by Theorems~\ref{thm:xpos} and~\ref{thm:slabs} and Proposition~\ref{prop:lower}.

\paragraph{Open problems.}
\begin{enumerate}
  \item \emph{The all-sector edge as a theorem.} Conjecture~\ref{conj:edge} for $t\ge3$. The ``if and only if'' on
  $(-1.5,x^*]$ in every sector: $j>12$ is open on the slices of Theorem~\ref{thm:slice}, $x\in(-0.02,0)$ is not covered by
  Theorem~\ref{thm:slabs}, and the small sectors $j<40$ on its boxes are not certified. The gauge-orthogonal edge in exact
  form.
  \item \emph{Near $t=2$.} The near-top families take over there ($c_+$ and $c_0$ vanish at $s=\sqrt2$). The all-sector edge
  near $t=2$ is unknown.
  \item \emph{The lower edge.} $x_{\mathrm{lo}}$ in closed form (only the bound $r_{\mathrm{lo}}$ is known), and the
  one-interval structure, which is computed on grids but not proved.
  \item \emph{Fields that mix sectors.} Circle-invariant fields that are not left-invariant, lifts by non-constant gauge
  transformations, other circles. There the twist is no longer alone in its charge block, so its exactness is itself open.
  \item \emph{Other groups.} In the $SU(3)$ family of Test~C, stretched along the hypercharge circle, the twist is
  the floor nowhere. Whether a family stretched along the circle of the twist's own charge behaves like $SU(2)$ is open,
  and would need its own pre-registration.
  \item \emph{Competitors in closed form.} The block $(3,1/2)$ has one edge factor of degrees $(11,14,26)$, irreducible
  over $\mathbb{Q}$. The polynomials of the small-$y$ blocks are known, but no two-dimensional region theorem at small $y$ is.
  \item \emph{A reading.} The decomposition of $\lambda$ invites reading the edge as the point where the field along the
  long circle turns against the twist's charge by enough, but not too much, as the shape of space sets. That is a
  hypothesis, not a result.
\end{enumerate}

\section*{Funding and competing interests}
This work received no funding. The author declares no competing interests.

\section*{Code and data availability}\label{sec:availability}
All programs and certificates are in the \textsc{GIGI} repository, \url{https://github.com/nurdymuny/gigi}
(tag \textbf{[TAG TO BE SUPPLIED]}). The exact parts use \textsc{SymPy}~\cite{sympy2017} and the interval parts
\textsc{mpmath}~\cite{mpmath}. Scripts by result:
\begin{itemize}
  \item the class of fields and the lower edge: \path{scripts/edge_wp4_design_checks.py};
  \item charge blocks and edge functions: \path{scripts/ricci_su2_exact_edge.py},
    \path{scripts/ricci_su2_charge_selection_symbolic.py};
  \item the closed form: \path{scripts/ricci_su2_twist_family_closed_form.py};
  \item the anatomy of the eigenvalue: \path{scripts/ricci_su2_twist_anatomy.py};
  \item the split of $R_0$: \path{scripts/ricci_su2_gyromagnetic_split.py};
  \item the edge polynomial: \path{scripts/ricci_su2_edge_polynomial.py};
  \item the competitor map and block polynomials: \path{scripts/edge_wp1_competitor_map.py},
    \path{scripts/edge_wp1_block_polynomials.py};
  \item Theorem~\ref{thm:xpos}: \path{scripts/edge_wp2_symbolic.py};
  \item Theorem~\ref{thm:slabs}: \path{scripts/edge_wp2_interval_proof.py};
  \item Theorem~\ref{thm:slice}: \path{scripts/edge_wp3_slice_iff.py};
  \item Theorem~\ref{thm:domain}: \path{scripts/ricci_su2_family_block_cert.py},
    \path{scripts/ricci_su2_family_large_j.py}, \path{scripts/ricci_su2_outward.py};
  \item the exploratory map: \path{scripts/ricci_su2_twist_family_map.py};
  \item the tests of Section~\ref{sec:decide}: \path{scripts/ricci_su2_test2_family.py},
    \path{scripts/ricci_su3_test2_family.py}, \path{scripts/ricci_alignment_test.py},
    \path{scripts/ricci_stretch_test.py}, \path{scripts/ricci_su2_edge_test.py}.
\end{itemize}
Receipts:
\begin{itemize}
  \item in \path{research/su2-floor-family-20261119/}: \path{twist_family_closed_form.txt}, \path{twist_anatomy.json},
    \path{gyromagnetic_split.txt}, \path{exact_edge.json}, \path{edge_polynomial.txt} and \path{.srepr},
    \path{family_block_cert_j80.json}, \path{family_large_j.json}, \path{charge_selection_symbolic.txt},
    \path{edge_test_*.json}, \path{alignment_test_*.json}, \path{stretch_test.json}, \path{test2_*.json},
    \path{su3_test2_*.json}, \path{twist_family_map.json};
  \item in \path{research/edge-theorem-20261121/}: \path{wp1_*}, \path{wp2_symbolic.txt},
    \path{wp2_interval_proof_*.json}, \path{wp2_large_sector_bound.md}, \path{wp3_slice_iff.md} and \path{.json},
    \path{wp4_design_checks.json}, \path{workflow_reports.json}, and the pre-registrations
    \path{wp4_prereg_A_heldout.md} (with \path{wp4_prereg_A_points.json}), \path{wp4_prereg_B_general_fields.md} and
    \path{wp4_prereg_C_su3.md}, with their amendments;
  \item the pre-registration notes of the earlier tests in \path{docs/research/}.
\end{itemize}
Sixteen algebraic statements, among them the factorisation of the edge determinant as $-16384\,P_1P_2$ and a crossing
of $P_2=0$ by the Berger path, are also checked in Lean~4~\cite{lean4} (v4.36.0) with Mathlib~\cite{mathlib2020}:
\path{docs/papers/berger-spectrum/lean/Helike/Price.lean} and \path{Helike/Edge.lean}, with \path{README.md} and
\path{RESULTS.md}. There are $16$ theorems, with no \texttt{sorry}, depending only on the axioms \texttt{propext},
\texttt{Classical.choice} and \texttt{Quot.sound}. They build with \texttt{lake exe cache get} and \texttt{lake build},
and \texttt{lake env lean Axioms.lean} lists the axioms. This paper, its source, and the code and data above are
archived on Zenodo (DOI \textbf{[TO BE ASSIGNED]}).

\begin{thebibliography}{99}\small
\bibitem{alefeldmayer1993} G.~Alefeld, G.~Mayer, The Cholesky method for interval data, \emph{Linear Algebra Appl.} 194 (1993) 161--182, \href{https://doi.org/10.1016/0024-3795(93)90120-D}{doi:10.1016/0024-3795(93)90120-D}.
\bibitem{avramidicollopy2012} I.~G.~Avramidi, S.~Collopy, Thermal Yang--Mills theory in the Einstein universe, \emph{J.\ Phys.\ A} 45 (2012) 374009, \href{https://doi.org/10.1088/1751-8113/45/37/374009}{doi:10.1088/1751-8113/45/37/374009}, \href{https://arxiv.org/abs/1201.5163}{arXiv:1201.5163}.
\bibitem{bar1992} C.~B\"ar, The Dirac operator on homogeneous spaces and its spectrum on $3$-dimensional lens spaces, \emph{Arch.\ Math.} 59 (1992) 65--79, \href{https://doi.org/10.1007/BF01199016}{doi:10.1007/BF01199016}.
\bibitem{bcgk2014} S.~Bolognesi, C.~Chatterjee, S.~B.~Gudnason, K.~Konishi, Vortex zero modes, large flux limit and Ambj{\o}rn--Nielsen--Olesen magnetic instabilities, \emph{JHEP} 10 (2014) 101, \href{https://doi.org/10.1007/JHEP10(2014)101}{doi:10.1007/JHEP10(2014)101}, \href{https://arxiv.org/abs/1408.1572}{arXiv:1408.1572}.
\bibitem{bl1981} J.-P.~Bourguignon, H.~B.~Lawson, Stability and isolation phenomena for Yang--Mills fields, \emph{Comm.\ Math.\ Phys.} 79 (1981) 189--230, \href{https://doi.org/10.1007/BF01942061}{doi:10.1007/BF01942061}.
\bibitem{davisfloor2026} B.~R.~Davis, The Floor Theorem for the Yang--Mills Hessian at the Berger field: a floor law across the Berger spheres, 2026, \url{https://doi.org/10.5281/zenodo.24412603} (concept DOI \href{https://doi.org/10.5281/zenodo.24386613}{10.5281/zenodo.24386613}).
\bibitem{haldane1983} F.~D.~M.~Haldane, Fractional quantization of the Hall effect: a hierarchy of incompressible quantum fluid states, \emph{Phys.\ Rev.\ Lett.} 51 (1983) 605--608, \href{https://doi.org/10.1103/PhysRevLett.51.605}{doi:10.1103/PhysRevLett.51.605}.
\bibitem{hanwen2026} X.~Han, Y.~Wen, Nonexistence of weakly stable Yang--Mills fields, 2026, \href{https://arxiv.org/abs/2601.17649}{arXiv:2601.17649}.
\bibitem{henkellauret2026} J.~Henkel, E.~A.~Lauret, Hodge Laplacian on $1$-forms of homogeneous $3$-spheres, 2026, \href{https://arxiv.org/abs/2605.05406}{arXiv:2605.05406}.
\bibitem{jin2026} X.~Jin, Stable abelian Yang--Mills--Higgs fields on the round three-sphere, 2026, \href{https://arxiv.org/abs/2609.24045}{arXiv:2609.24045}.
\bibitem{klc2021} K.~Kumar, O.~Lechtenfeld, G.~Pican\c{c}o Costa, Instability of cosmic Yang--Mills fields, \emph{Nucl.\ Phys.\ B} 973 (2021) 115583, \href{https://doi.org/10.1016/j.nuclphysb.2021.115583}{doi:10.1016/j.nuclphysb.2021.115583}, \href{https://arxiv.org/abs/2102.08401}{arXiv:2102.08401}.
\bibitem{klingschueth2022} J.~Kling, D.~Schueth, On the Dirac spectrum of homogeneous $3$-spheres, \emph{J.\ Geom.\ Anal.} 32 (2022), no.~11, Paper No.~175, \href{https://doi.org/10.1007/s12220-022-00997-x}{doi:10.1007/s12220-022-00997-x}, \href{https://arxiv.org/abs/2204.12990}{arXiv:2204.12990}.
\bibitem{kot1986} S.~Kobayashi, Y.~Ohnita, M.~Takeuchi, On instability of Yang--Mills connections, \emph{Math.\ Z.} 193 (1986) 165--189, \href{https://doi.org/10.1007/BF01174329}{doi:10.1007/BF01174329}.
\bibitem{laquer1984} H.~T.~Laquer, Stability properties of the Yang--Mills functional near the canonical connection, \emph{Michigan Math.\ J.} 31 (1984) 139--159, \href{https://doi.org/10.1307/mmj/1029003019}{doi:10.1307/mmj/1029003019}.
\bibitem{lauret2019} E.~A.~Lauret, The smallest Laplace eigenvalue of homogeneous $3$-spheres, \emph{Bull.\ Lond.\ Math.\ Soc.} 51 (2019) 49--69, \href{https://doi.org/10.1112/blms.12213}{doi:10.1112/blms.12213}, \href{https://arxiv.org/abs/1801.04259}{arXiv:1801.04259}.
\bibitem{lean4} L.~de Moura, S.~Ullrich, The Lean 4 theorem prover and programming language, in \emph{Automated Deduction -- CADE 28}, Lecture Notes in Comput.\ Sci.\ 12699, Springer, 2021, 625--635, \href{https://doi.org/10.1007/978-3-030-79876-5_37}{doi:10.1007/978-3-030-79876-5\_37}.
\bibitem{mathlib2020} The mathlib Community, The Lean mathematical library, in \emph{Proc.\ 9th ACM SIGPLAN Int.\ Conf.\ Certified Programs and Proofs (CPP 2020)}, 2020, 367--381, \href{https://doi.org/10.1145/3372885.3373824}{doi:10.1145/3372885.3373824}.
\bibitem{mpmath} The mpmath development team, mpmath: a Python library for arbitrary-precision floating-point arithmetic, \url{https://mpmath.org}.
\bibitem{nielsenolesen1978} N.~K.~Nielsen, P.~Olesen, An unstable Yang--Mills field mode, \emph{Nucl.\ Phys.\ B} 144 (1978) 376--396, \href{https://doi.org/10.1016/0550-3213(78)90377-2}{doi:10.1016/0550-3213(78)90377-2}.
\bibitem{ohnitapan1990} Y.~Ohnita, Y.~L.~Pan, On weakly stable Yang--Mills fields over positively pinched manifolds and certain symmetric spaces, \emph{Kodai Math.\ J.} 13 (1990) 317--332, \href{https://doi.org/10.2996/kmj/1138039277}{doi:10.2996/kmj/1138039277}.
\bibitem{savvidy1977} G.~K.~Savvidy, Infrared instability of the vacuum state of gauge theories and asymptotic freedom, \emph{Phys.\ Lett.\ B} 71 (1977) 133--134, \href{https://doi.org/10.1016/0370-2693(77)90759-6}{doi:10.1016/0370-2693(77)90759-6}.
\bibitem{stern2010} M.~Stern, Geometry of minimal energy Yang--Mills connections, \emph{J.\ Differential Geom.} 86 (2010) 163--188, \href{https://doi.org/10.4310/jdg/1299766686}{doi:10.4310/jdg/1299766686}, \href{https://arxiv.org/abs/0808.0667}{arXiv:0808.0667}.
\bibitem{sympy2017} A.~Meurer et al., SymPy: symbolic computing in Python, \emph{PeerJ Comput.\ Sci.} 3 (2017) e103, \href{https://doi.org/10.7717/peerj-cs.103}{doi:10.7717/peerj-cs.103}.
\bibitem{tanjiitakura2012} N.~Tanji, K.~Itakura, Schwinger mechanism enhanced by the Nielsen--Olesen instability, \emph{Phys.\ Lett.\ B} (2012), \href{https://doi.org/10.1016/j.physletb.2012.05.043}{doi:10.1016/j.physletb.2012.05.043}, \href{https://arxiv.org/abs/1111.6772}{arXiv:1111.6772}.
\bibitem{varga2001} R.~S.~Varga, Ger\v{s}gorin-type eigenvalue inclusion theorems and their sharpness, \emph{Electron.\ Trans.\ Numer.\ Anal.} 12 (2001) 113--133.
\bibitem{vonneumannwigner1929} J.~von Neumann, E.~Wigner, \"Uber das Verhalten von Eigenwerten bei adiabatischen Prozessen, \emph{Phys.\ Z.} 30 (1929) 467--470.
\bibitem{wuyang1976} T.~T.~Wu, C.~N.~Yang, Dirac monopole without strings: monopole harmonics, \emph{Nucl.\ Phys.\ B} 107 (1976) 365--380, \href{https://doi.org/10.1016/0550-3213(76)90143-7}{doi:10.1016/0550-3213(76)90143-7}.
\end{thebibliography}
\end{document}
